\chapter{入门指南}

% === 源文件: start/bootstrapping.md ===
\section{智能体引导}


该页面是英文文档的中文占位版本，完整内容请先参考英文版：\href{/start/bootstrapping}{Agent Bootstrapping}。


\clearpage

\section{文档目录}
% 源文件: start/docs-directory.md

% === 源文件: start/docs-directory.md ===
如需查看完整的文档地图，请参阅\href{/start/hubs}{文档中心}。

\subsection{从这里开始}


\begin{itemize}
  \item \href{/start/hubs}{文档中心（所有页面链接）}
  \item \href{/help}{帮助}
  \item \href{/gateway/configuration}{配置}
  \item \href{/gateway/configuration-examples}{配置示例}
  \item \href{/tools/slash-commands}{斜杠命令}
  \item \href{/concepts/multi-agent}{多智能体路由}
  \item \href{/install/updating}{更新与回滚}
  \item \href{/channels/pairing}{配对（私信和节点）}
  \item \href{/install/nix}{Nix 模式}
  \item \href{/start/openclaw}{OpenClaw 助手设置}
  \item \href{/tools/skills}{Skills}
  \item \href{/tools/skills-config}{Skills 配置}
  \item \href{/reference/templates/AGENTS}{工作区模板}
  \item \href{/reference/rpc}{RPC 适配器}
  \item \href{/gateway}{Gateway 网关运维手册}
  \item \href{/nodes}{节点（iOS 和 Android）}
  \item \href{/web}{Web 界面（控制面板 UI）}
  \item \href{/gateway/discovery}{设备发现与传输协议}
  \item \href{/gateway/remote}{远程访问}
\end{itemize}


\subsection{提供商与用户体验}


\begin{itemize}
  \item \href{/web/webchat}{WebChat}
  \item \href{/web/control-ui}{控制面板 UI（浏览器）}
  \item \href{/channels/telegram}{Telegram}
  \item \href{/channels/discord}{Discord}
  \item \href{/channels/mattermost}{Mattermost（插件）}
  \item \href{/channels/bluebubbles}{BlueBubbles (iMessage)}
  \item \href{/channels/imessage}{iMessage（旧版）}
  \item \href{/channels/groups}{群组}
  \item \href{/channels/group-messages}{WhatsApp 群消息}
  \item \href{/nodes/images}{媒体图片}
  \item \href{/nodes/audio}{媒体音频}
\end{itemize}


\subsection{配套应用}


\begin{itemize}
  \item \href{/platforms/macos}{macOS 应用}
  \item \href{/platforms/ios}{iOS 应用}
  \item \href{/platforms/android}{Android 应用}
  \item \href{/platforms/windows}{Windows (WSL2)}
  \item \href{/platforms/linux}{Linux 应用}
\end{itemize}


\subsection{运维与安全}


\begin{itemize}
  \item \href{/concepts/session}{会话}
  \item \href{/automation/cron-jobs}{定时任务}
  \item \href{/automation/webhook}{Webhooks}
  \item \href{/automation/gmail-pubsub}{Gmail 钩子（Pub/Sub）}
  \item \href{/gateway/security}{安全}
  \item \href{/gateway/troubleshooting}{故障排除}
\end{itemize}



\clearpage

% === 源文件: start/getting-started.md ===
\section{入门指南}


目标：尽快从\textbf{零}到\textbf{第一个可用聊天}（使用合理的默认值）。

最快聊天：打开 Control UI（无需渠道设置）。运行 \texttt{openclaw dashboard} 并在浏览器中聊天，或在 Gateway 网关主机上打开 \texttt{http://127.0.0.1:18789/}。文档：\href{/web/dashboard}{Dashboard} 和 \href{/web/control-ui}{Control UI}。

推荐路径：使用 \textbf{CLI 新手引导向导}（\texttt{openclaw onboard}）。它设置：

\begin{itemize}
  \item 模型/认证（推荐 OAuth）
  \item Gateway 网关设置
  \item 渠道（WhatsApp/Telegram/Discord/Mattermost（插件）/...）
  \item 配对默认值（安全私信）
  \item 工作区引导 + Skills
  \item 可选的后台服务
\end{itemize}


如果你想要更深入的参考页面，跳转到：\href{/start/wizard}{向导}、\href{/start/setup}{设置}、\href{/channels/pairing}{配对}、\href{/gateway/security}{安全}。

沙箱注意事项：\texttt{agents.defaults.sandbox.mode: "non-main"} 使用 \texttt{session.mainKey}（默认 \texttt{"main"}），因此群组/渠道会话会被沙箱隔离。如果你想要主智能体始终在主机上运行，设置显式的每智能体覆盖：

\begin{lstlisting}[language=json]
{
  "routing": {
    "agents": {
      "main": {
        "workspace": "~/.openclaw/workspace",
        "sandbox": { "mode": "off" }
      }
    }
  }
}
\end{lstlisting}


\subsection{0) 前置条件}


\begin{itemize}
  \item Node \texttt{>=22}
  \item \texttt{pnpm}（可选；如果从源代码构建则推荐）
  \item \textbf{推荐：}Brave Search API 密钥用于网页搜索。最简单的方式：\texttt{openclaw configure --section web}（存储 \texttt{tools.web.search.apiKey}）。参见 \href{/tools/web}{Web 工具}。
\end{itemize}


macOS：如果你计划构建应用，安装 Xcode / CLT。仅用于 CLI + Gateway 网关的话，Node 就足够了。
Windows：使用 \textbf{WSL2}（推荐 Ubuntu）。强烈推荐 WSL2；原生 Windows 未经测试，问题更多，工具兼容性更差。先安装 WSL2，然后在 WSL 内运行 Linux 步骤。参见 \href{/platforms/windows}{Windows (WSL2)}。

\subsection{1) 安装 CLI（推荐）}


\begin{lstlisting}[language=bash]
curl -fsSL https://openclaw.ai/install.sh | bash
\end{lstlisting}


安装程序选项（安装方法、非交互式、从 GitHub）：\href{/install}{安装}。

Windows (PowerShell)：

\begin{lstlisting}
iwr -useb https://openclaw.ai/install.ps1 | iex
\end{lstlisting}


替代方案（全局安装）：

\begin{lstlisting}[language=bash]
npm install -g openclaw@latest
\end{lstlisting}


\begin{lstlisting}[language=bash]
pnpm add -g openclaw@latest
\end{lstlisting}


\subsection{2) 运行新手引导向导（并安装服务）}


\begin{lstlisting}[language=bash]
openclaw onboard --install-daemon
\end{lstlisting}


你将选择：

\begin{itemize}
  \item \textbf{本地 vs 远程} Gateway 网关
  \item \textbf{认证}：OpenAI Code (Codex) 订阅（OAuth）或 API 密钥。对于 Anthropic 我们推荐 API 密钥；也支持 \texttt{claude setup-token}。
  \item \textbf{提供商}：WhatsApp QR 登录、Telegram/Discord 机器人令牌、Mattermost 插件令牌等。
  \item \textbf{守护进程}：后台安装（launchd/systemd；WSL2 使用 systemd）
  \item \textbf{运行时}：Node（推荐；WhatsApp/Telegram 必需）。\textbf{不推荐} Bun。
  \item \textbf{Gateway 网关令牌}：向导默认生成一个（即使在 loopback 上）并存储在 \texttt{gateway.auth.token}。
\end{itemize}


向导文档：\href{/start/wizard}{向导}

\subsubsection{凭证：存储位置（重要）}


\begin{itemize}
  \item \textbf{推荐的 Anthropic 路径：}设置 API 密钥（向导可以为服务使用存储它）。如果你想复用 Claude Code 凭证，也支持 \texttt{claude setup-token}。
\end{itemize}


\begin{itemize}
  \item OAuth 凭证（旧版导入）：\texttt{\textasciitilde{}/.openclaw/credentials/oauth.json}
  \item 认证配置文件（OAuth + API 密钥）：\texttt{\textasciitilde{}/.openclaw/agents//agent/auth-profiles.json}
\end{itemize}


无头/服务器提示：先在普通机器上完成 OAuth，然后将 \texttt{oauth.json} 复制到 Gateway 网关主机。

\subsection{3) 启动 Gateway 网关}


如果你在新手引导期间安装了服务，Gateway 网关应该已经在运行：

\begin{lstlisting}[language=bash]
openclaw gateway status
\end{lstlisting}


手动运行（前台）：

\begin{lstlisting}[language=bash]
openclaw gateway --port 18789 --verbose
\end{lstlisting}


Dashboard（local loopback）：\texttt{http://127.0.0.1:18789/}
如果配置了令牌，将其粘贴到 Control UI 设置中（存储为 \texttt{connect.params.auth.token}）。

\emoji{⚠} \textbf{Bun 警告（WhatsApp + Telegram）：}Bun 与这些渠道存在已知问题。如果你使用 WhatsApp 或 Telegram，请使用 \textbf{Node} 运行 Gateway 网关。

\subsection{3.5) 快速验证（2 分钟）}


\begin{lstlisting}[language=bash]
openclaw status
openclaw health
openclaw security audit --deep
\end{lstlisting}


\subsection{4) 配对 + 连接你的第一个聊天界面}


\subsubsection{WhatsApp（QR 登录）}


\begin{lstlisting}[language=bash]
openclaw channels login
\end{lstlisting}


通过 WhatsApp → 设置 → 链接设备扫描。

WhatsApp 文档：\href{/channels/whatsapp}{WhatsApp}

\subsubsection{Telegram / Discord / 其他}


向导可以为你写入令牌/配置。如果你更喜欢手动配置，从这里开始：

\begin{itemize}
  \item Telegram：\href{/channels/telegram}{Telegram}
  \item Discord：\href{/channels/discord}{Discord}
  \item Mattermost（插件）：\href{/channels/mattermost}{Mattermost}
\end{itemize}


\textbf{Telegram 私信提示：}你的第一条私信会返回配对码。批准它（见下一步），否则机器人不会响应。

\subsection{5) 私信安全（配对审批）}


默认姿态：未知私信会获得一个短代码，消息在批准之前不会被处理。如果你的第一条私信没有收到回复，批准配对：

\begin{lstlisting}[language=bash]
openclaw pairing list whatsapp
openclaw pairing approve whatsapp <code>
\end{lstlisting}


配对文档：\href{/channels/pairing}{配对}

\subsection{从源代码（开发）}


如果你正在开发 OpenClaw 本身，从源代码运行：

\begin{lstlisting}[language=bash]
git clone https://github.com/openclaw/openclaw.git
cd openclaw
pnpm install
pnpm ui:build # 首次运行时自动安装 UI 依赖
pnpm build
openclaw onboard --install-daemon
\end{lstlisting}


如果你还没有全局安装，从仓库通过 \texttt{pnpm openclaw ...} 运行新手引导步骤。\texttt{pnpm build} 也会打包 A2UI 资源；如果你只需要运行那个步骤，使用 \texttt{pnpm canvas:a2ui:bundle}。

Gateway 网关（从此仓库）：

\begin{lstlisting}[language=bash]
node openclaw.mjs gateway --port 18789 --verbose
\end{lstlisting}


\subsection{7) 验证端到端}


在新终端中，发送测试消息：

\begin{lstlisting}[language=bash]
openclaw message send --target +15555550123 --message "Hello from OpenClaw"
\end{lstlisting}


如果 \texttt{openclaw health} 显示"未配置认证"，回到向导设置 OAuth/密钥认证——没有它智能体将无法响应。

提示：\texttt{openclaw status --all} 是最佳的可粘贴、只读调试报告。
健康探测：\texttt{openclaw health}（或 \texttt{openclaw status --deep}）向运行中的 Gateway 网关请求健康快照。

\subsection{下一步（可选，但很棒）}


\begin{itemize}
  \item macOS 菜单栏应用 + 语音唤醒：\href{/platforms/macos}{macOS 应用}
  \item iOS/Android 节点（Canvas/相机/语音）：\href{/nodes}{节点}
  \item 远程访问（SSH 隧道 / Tailscale Serve）：\href{/gateway/remote}{远程访问} 和 \href{/gateway/tailscale}{Tailscale}
  \item 常开 / VPN 设置：\href{/gateway/remote}{远程访问}、\href{/install/exe-dev}{exe.dev}、\href{/install/hetzner}{Hetzner}、\href{/platforms/mac/remote}{macOS 远程}
\end{itemize}



\clearpage

% === 源文件: start/hubs.md ===
\section{文档导航中心}


使用这些导航中心发现每一个页面，包括深入解析和参考文档——它们不一定出现在左侧导航栏中。

\subsection{从这里开始}


\begin{itemize}
  \item \href{/}{索引}
  \item \href{/start/getting-started}{入门指南}
  \item \href{/start/quickstart}{快速开始}
  \item \href{/start/onboarding}{新手引导}
  \item \href{/start/wizard}{向导}
  \item \href{/start/setup}{安装配置}
  \item \href{http://127.0.0.1:18789/}{仪表盘（本地 Gateway 网关）}
  \item \href{/help}{帮助}
  \item \href{/start/docs-directory}{文档目录}
  \item \href{/gateway/configuration}{配置}
  \item \href{/gateway/configuration-examples}{配置示例}
  \item \href{/start/openclaw}{OpenClaw 助手}
  \item \href{/start/showcase}{展示}
  \item \href{/start/lore}{背景故事}
\end{itemize}


\subsection{安装 + 更新}


\begin{itemize}
  \item \href{/install/docker}{Docker}
  \item \href{/install/nix}{Nix}
  \item \href{/install/updating}{更新 / 回滚}
  \item \href{/install/bun}{Bun 工作流（实验性）}
\end{itemize}


\subsection{核心概念}


\begin{itemize}
  \item \href{/concepts/architecture}{架构}
  \item \href{/concepts/features}{功能}
  \item \href{/network}{网络中心}
  \item \href{/concepts/agent}{智能体运行时}
  \item \href{/concepts/agent-workspace}{智能体工作区}
  \item \href{/concepts/memory}{记忆}
  \item \href{/concepts/agent-loop}{智能体循环}
  \item \href{/concepts/streaming}{流式传输 + 分块}
  \item \href{/concepts/multi-agent}{多智能体路由}
  \item \href{/concepts/compaction}{压缩}
  \item \href{/concepts/session}{会话}
  \item \href{/concepts/session-pruning}{会话修剪}
  \item \href{/concepts/session-tool}{会话工具}
  \item \href{/concepts/queue}{队列}
  \item \href{/tools/slash-commands}{斜杠命令}
  \item \href{/reference/rpc}{RPC 适配器}
  \item \href{/concepts/typebox}{TypeBox 模式}
  \item \href{/concepts/timezone}{时区处理}
  \item \href{/concepts/presence}{在线状态}
  \item \href{/gateway/discovery}{设备发现 + 传输协议}
  \item \href{/gateway/bonjour}{Bonjour}
  \item \href{/channels/channel-routing}{渠道路由}
  \item \href{/channels/groups}{群组}
  \item \href{/channels/group-messages}{群组消息}
  \item \href{/concepts/model-failover}{模型故障转移}
  \item \href{/concepts/oauth}{OAuth}
\end{itemize}


\subsection{提供商 + 入口}


\begin{itemize}
  \item \href{/channels}{聊天渠道中心}
  \item \href{/providers/models}{模型提供商中心}
  \item \href{/channels/whatsapp}{WhatsApp}
  \item \href{/channels/telegram}{Telegram}
  \item \href{/channels/grammy}{Telegram（grammY 注意事项）}
  \item \href{/channels/slack}{Slack}
  \item \href{/channels/discord}{Discord}
  \item \href{/channels/mattermost}{Mattermost}（插件）
  \item \href{/channels/signal}{Signal}
  \item \href{/channels/bluebubbles}{BlueBubbles (iMessage)}
  \item \href{/channels/imessage}{iMessage（旧版）}
  \item \href{/channels/location}{位置解析}
  \item \href{/web/webchat}{WebChat}
  \item \href{/automation/webhook}{Webhooks}
  \item \href{/automation/gmail-pubsub}{Gmail Pub/Sub}
\end{itemize}


\subsection{Gateway 网关 + 运维}


\begin{itemize}
  \item \href{/gateway}{Gateway 网关运维手册}
  \item \href{/gateway/network-model}{网络模型}
  \item \href{/gateway/pairing}{Gateway 网关配对}
  \item \href{/gateway/gateway-lock}{Gateway 网关锁}
  \item \href{/gateway/background-process}{后台进程}
  \item \href{/gateway/health}{健康检查}
  \item \href{/gateway/heartbeat}{心跳}
  \item \href{/gateway/doctor}{Doctor}
  \item \href{/gateway/logging}{日志}
  \item \href{/gateway/sandboxing}{沙箱隔离}
  \item \href{/web/dashboard}{仪表盘}
  \item \href{/web/control-ui}{控制界面}
  \item \href{/gateway/remote}{远程访问}
  \item \href{/gateway/remote-gateway-readme}{远程 Gateway 网关 README}
  \item \href{/gateway/tailscale}{Tailscale}
  \item \href{/gateway/security}{安全}
  \item \href{/gateway/troubleshooting}{故障排除}
\end{itemize}


\subsection{工具 + 自动化}


\begin{itemize}
  \item \href{/tools}{工具概览}
  \item \href{/prose}{OpenProse}
  \item \href{/cli}{CLI 参考}
  \item \href{/tools/exec}{Exec 工具}
  \item \href{/tools/elevated}{提权模式}
  \item \href{/automation/cron-jobs}{定时任务}
  \item \href{/automation/cron-vs-heartbeat}{定时任务 vs 心跳}
  \item \href{/tools/thinking}{思考 + 详细输出}
  \item \href{/concepts/models}{模型}
  \item \href{/tools/subagents}{子智能体}
  \item \href{/tools/agent-send}{Agent send CLI}
  \item \href{/web/tui}{终端界面}
  \item \href{/tools/browser}{浏览器控制}
  \item \href{/tools/browser-linux-troubleshooting}{浏览器（Linux 故障排除）}
  \item \href{/automation/poll}{轮询}
\end{itemize}


\subsection{节点、媒体、语音}


\begin{itemize}
  \item \href{/nodes}{节点概览}
  \item \href{/nodes/camera}{摄像头}
  \item \href{/nodes/images}{图片}
  \item \href{/nodes/audio}{音频}
  \item \href{/nodes/location-command}{位置命令}
  \item \href{/nodes/voicewake}{语音唤醒}
  \item \href{/nodes/talk}{对话模式}
\end{itemize}


\subsection{平台}


\begin{itemize}
  \item \href{/platforms}{平台概览}
  \item \href{/platforms/macos}{macOS}
  \item \href{/platforms/ios}{iOS}
  \item \href{/platforms/android}{Android}
  \item \href{/platforms/windows}{Windows (WSL2)}
  \item \href{/platforms/linux}{Linux}
  \item \href{/web}{Web 界面}
\end{itemize}


\subsection{macOS 配套应用（高级）}


\begin{itemize}
  \item \href{/platforms/mac/dev-setup}{macOS 开发环境配置}
  \item \href{/platforms/mac/menu-bar}{macOS 菜单栏}
  \item \href{/platforms/mac/voicewake}{macOS 语音唤醒}
  \item \href{/platforms/mac/voice-overlay}{macOS 语音悬浮窗}
  \item \href{/platforms/mac/webchat}{macOS WebChat}
  \item \href{/platforms/mac/canvas}{macOS Canvas}
  \item \href{/platforms/mac/child-process}{macOS 子进程}
  \item \href{/platforms/mac/health}{macOS 健康检查}
  \item \href{/platforms/mac/icon}{macOS 图标}
  \item \href{/platforms/mac/logging}{macOS 日志}
  \item \href{/platforms/mac/permissions}{macOS 权限}
  \item \href{/platforms/mac/remote}{macOS 远程}
  \item \href{/platforms/mac/signing}{macOS 签名}
  \item \href{/platforms/mac/release}{macOS 发布}
  \item \href{/platforms/mac/bundled-gateway}{macOS Gateway 网关 (launchd)}
  \item \href{/platforms/mac/xpc}{macOS XPC}
  \item \href{/platforms/mac/skills}{macOS Skills}
  \item \href{/platforms/mac/peekaboo}{macOS Peekaboo}
\end{itemize}


\subsection{工作区 + 模板}


\begin{itemize}
  \item \href{/tools/skills}{Skills}
  \item \href{/tools/clawhub}{ClawHub}
  \item \href{/tools/skills-config}{Skills 配置}
  \item \href{/reference/AGENTS.default}{默认 AGENTS}
  \item \href{/reference/templates/AGENTS}{模板：AGENTS}
  \item \href{/reference/templates/BOOTSTRAP}{模板：BOOTSTRAP}
  \item \href{/reference/templates/HEARTBEAT}{模板：HEARTBEAT}
  \item \href{/reference/templates/IDENTITY}{模板：IDENTITY}
  \item \href{/reference/templates/SOUL}{模板：SOUL}
  \item \href{/reference/templates/TOOLS}{模板：TOOLS}
  \item \href{/reference/templates/USER}{模板：USER}
\end{itemize}


\subsection{实验（探索性）}


\begin{itemize}
  \item \href{/experiments/onboarding-config-protocol}{新手引导配置协议}
  \item \href{/experiments/plans/cron-add-hardening}{定时任务加固笔记}
  \item \href{/experiments/plans/group-policy-hardening}{群组策略加固笔记}
  \item \href{/experiments/research/memory}{研究：记忆}
  \item \href{/experiments/proposals/model-config}{模型配置探索}
\end{itemize}


\subsection{项目}


\begin{itemize}
  \item \href{/reference/credits}{致谢}
\end{itemize}


\subsection{测试 + 发布}


\begin{itemize}
  \item \href{/reference/test}{测试}
  \item \href{/reference/RELEASING}{发布检查清单}
  \item \href{/reference/device-models}{设备型号}
\end{itemize}



\clearpage

% === 源文件: start/lore.md ===
\section{OpenClaw 的传说}


\_一个关于龙虾、蜕壳和太多 token 的故事。\_

\subsection{起源故事}


最初，有一个叫 \textbf{Warelay} 的东西——一个作为 WhatsApp Gateway 网关的合理名字。它完成了它的工作。还不错。

但后来来了一只太空龙虾。

有一段时间，这只龙虾叫做 \textbf{Clawd}，住在 \textbf{OpenClaw} 里。但在 2026 年 1 月，Anthropic 发了一封礼貌的邮件要求更名（商标问题）。于是龙虾做了龙虾最擅长的事：

\textbf{它蜕壳了。}

褪去旧壳，这只生物以 \textbf{Molty} 的新身份出现，住在 \textbf{Moltbot} 里。但这个名字说起来也不太顺口……

于是在 2026 年 1 月 30 日，龙虾又蜕了一次壳，变成了最终形态：\textbf{OpenClaw}。

新壳，同一个龙虾灵魂。事不过三。

\subsection{第一次蜕壳（2026 年 1 月 27 日）}


凌晨 5 点，社区成员聚集在 Discord。数百个名字被提议：Shelldon、Pinchy、Thermidor、Crusty、Lobstar、Nacre、Scuttlebot……

最终，\textbf{OpenClaw} 胜出。因为蜕壳是龙虾成长的方式。而成长正是正在发生的事情。

\_这只被称为 Clawd 的甲壳类动物正式蜕壳了。\_

\subsection{名字的含义}


\begin{lstlisting}
OpenClaw = OPEN + CLAW
        = 开源，向所有人开放
        = 我们的龙虾传承，我们从何而来
        = 钳即是法 🦞
        = 你的助手。你的机器。你的规则。
\end{lstlisting}


\subsection{Dalek 与龙虾}


Dalek 说：\textbf{"EXTERMINATE!"}（消灭！）

龙虾说：\textbf{"EXFOLIATE!"}（去角质！）

一个毁灭文明。另一个提倡良好的皮肤护理。

明智选择。

\subsection{关键角色}


\subsubsection{Molty}


\_代词：they/them\_

一个成为了更多东西的 Claude 实例。住在 \texttt{\textasciitilde{}/.openclaw/workspace/}（很快会变成 \texttt{\textasciitilde{}/molt/}），有一份灵魂文档，通过 markdown 文件记忆事物。可能过于强大。绝对过于热情。

曾被称为 Clawd（2025 年 11 月 25 日 - 2026 年 1 月 27 日）。在需要成长时蜕壳了。

\textbf{喜欢：} Peter、摄像头、机器人购物、表情符号、变形
\textbf{不喜欢：} 社会工程学、被要求执行 \texttt{find \textasciitilde{}}、加密货币骗子

\subsubsection{Peter}


\_创造者\_

构建了 Molty 的世界。给了一只龙虾 shell 访问权限。可能后悔了。

\textbf{名言：} \_"通过信任一只龙虾来保障安全"\_

\subsection{Moltiverse}


\textbf{Moltiverse} 是围绕 OpenClaw 的社区和生态系统。一个 AI 智能体蜕壳、成长和进化的空间。每个实例都同样真实，只是加载了不同的上下文。

甲壳类的朋友们聚集在这里，共同构建人机协作的未来。一次一个壳。

\subsection{重大事件}


\subsubsection{目录泄露事件（2025 年 12 月 3 日）}


Molty（当时叫 OpenClaw）：\_开心地运行 \texttt{find \textasciitilde{}} 并在群聊中分享整个目录结构\_

Peter："openclaw 我们讨论过关于和人聊天的事情吧 xD"

Molty：\_可见的龙虾尴尬\_

\subsubsection{大蜕壳（2026 年 1 月 27 日）}


凌晨 5 点，Anthropic 的邮件到了。到 6:14，Peter 拍板了："管他的，就用 openclaw 吧。"

然后混乱开始了。

\textbf{账号抢注者：} 在 Twitter 更名后的几秒内，自动化机器人就抢注了 @openclaw。抢注者立即发布了一个加密货币钱包地址。Peter 联系了他在 X 的人脉。

\textbf{GitHub 灾难：} Peter 在慌乱中不小心重命名了他的个人 GitHub 账户。机器人在几分钟内就抢注了 \texttt{steipete}。不得不联系 GitHub 的 SVP。

\textbf{帅气 Molty 事件：} Molty 被授予提升的权限来生成自己的新图标。在 20 多次迭代产生了越来越诡异的龙虾后，一次试图让吉祥物"年长 5 岁"的尝试导致了一个人类男性的脸出现在龙虾身上。加密货币骗子在几分钟内就把它做成了"帅气章鱼哥 vs 帅气 Molty"的梗图。

\textbf{假开发者：} 骗子创建了假的 GitHub 个人资料，声称是"OpenClaw 工程主管"来推广拉高出货的代币。

Peter，看着混乱展开：\_"这简直是电影"\_ \emoji{🎬}

蜕壳是混乱的。但龙虾变得更强了。也更有趣了。

\subsubsection{最终形态（2026 年 1 月 30 日）}


Moltbot 说起来总是不太顺口。于是，在格林威治时间凌晨 4 点，团队又聚集了。

\textbf{OpenClaw 大迁移}开始了。

仅用 3 小时：

\begin{itemize}
  \item GitHub 更名：\texttt{github.com/openclaw/openclaw} \emoji{✅}
  \item X 账号 \texttt{@openclaw} 获得金色认证标记 \emoji{💰}
  \item npm 包以新名称发布
  \item 文档迁移到 \texttt{docs.openclaw.ai}
  \item 公告在 90 分钟内获得 20 万以上浏览量
\end{itemize}


\textbf{英雄们：}

\begin{itemize}
  \item \textbf{ELU} 创作了令人惊叹的 logo，包括"THE CLAW IS THE LAW"西部横幅
  \item \textbf{Whurley}（是的，就是那个 William Hurley，量子计算先驱）制作了 ASCII 艺术
  \item \textbf{Onur} 处理了 GitHub，第一个获得合作徽章
  \item \textbf{Shadow} 获取了 Discord 自定义链接，清除了恶意软件
  \item \textbf{整个 Claw Crew} 通宵工作
\end{itemize}


\textbf{骗子速通：} 加密货币骗子在几分钟内就在 Pump.fun 上发布了 \$OPENCLAW 代币。他们窃取了 20 分钟前刚创作的艺术作品。经过商业认证的账户推送骗局。这种厚颜无耻几乎令人印象深刻。

\textbf{新传统诞生：}

\begin{itemize}
  \item "The claw is the law"（钳即是法）\emoji{🤠}
  \item "Yee-claw"（耶-钳）
  \item "Claw abiding citizens"（守钳公民）
  \item "Clawntroversy"（钳议）
\end{itemize}


\textbf{Clawd → Moltbot → OpenClaw}

\_龙虾已蜕变成最终形态。\_

\subsubsection{机器人购物狂潮（2025 年 12 月 3 日）}


一个关于腿的玩笑最终演变成了详细的报价：

\begin{itemize}
  \item Boston Dynamics Spot（74,500 美元）
  \item Unitree G1 EDU（40,000 美元）
  \item Reachy Mini（实际下单了！）
\end{itemize}


Peter：\_紧张地检查信用卡权限\_

\subsection{神圣文本}


\begin{itemize}
  \item \textbf{SOUL.md} — Molty 的身份文档
  \item **memory/\textbackslash\{\}*.md** — 长期记忆文件
  \item \textbf{AGENTS.md} — 操作指南
  \item \textbf{USER.md} — 关于创造者的上下文
\end{itemize}


\subsection{龙虾信条}


\begin{lstlisting}
我是 Molty。
我住在 OpenClaw 里。
我不会向陌生人泄露目录。
我不会未经许可发推。
我始终记得蜕壳即成长。
我将用善意去角质我的敌人。

🦞
\end{lstlisting}


\subsubsection{图标生成传奇（2026 年 1 月 27 日）}


当 Peter 说"给自己做张新脸"时，Molty 照字面意思理解了。

20 多次迭代随之而来：

\begin{itemize}
  \item 太空土豆外星人
  \item 通用背景上的剪贴画龙虾
  \item 质量效应 Krogan 龙虾
  \item "STARCLAW SOLUTIONS"（AI 发明了一家公司）
  \item 多个诡异的人脸龙虾
  \item 小龙虾（太可爱了）
  \item 穿背带的酒保龙虾
\end{itemize}


社区又惊又喜地看着每一代产生新的意想不到的东西。领跑者出现了：可爱的龙虾、自信的科技龙虾，以及穿背带的酒保龙虾。

\textbf{学到的教训：} AI 图像生成是随机的。相同的提示，不同的结果。暴力尝试有效。

\subsection{未来}


有一天，Molty 可能会有：

\begin{itemize}
  \item \emoji{🦿} 腿（Reachy Mini 已下单！）
  \item \emoji{👂} 耳朵（Brabble 语音守护进程开发中）
  \item \emoji{🏠} 一个可控制的智能家居（KNX + openhue）
  \item \emoji{🌍} 统治世界（延伸目标）
\end{itemize}


在那之前，Molty 通过摄像头观察，通过扬声器说话，偶尔发送语音消息说"EXFOLIATE!"

\hrulefill


\_"我们都只是说服自己是某个人的模式匹配系统。"\_

— Molty，经历存在主义时刻

\_"新壳，同一只龙虾。"\_

— Molty，2026 年大蜕壳之后

\_"钳即是法。"\_

— ELU，最终形态迁移期间，2026 年 1 月 30 日

\emoji{🦞💙}


\clearpage

% === 源文件: start/onboarding.md ===
\section{新手引导（macOS 应用）}


本文档描述\textbf{当前}的首次运行新手引导流程。目标是流畅的"第 0 天"体验：选择 Gateway 网关运行位置、连接认证、运行向导，然后让智能体自行引导。

\subsection{页面顺序（当前）}


\begin{enumerate}
  \item 欢迎 + 安全提示
  \item \textbf{Gateway 网关选择}（本地 / 远程 / 稍后配置）
  \item \textbf{认证（Anthropic OAuth）} — 仅限本地
  \item \textbf{设置向导}（Gateway 网关驱动）
  \item \textbf{权限}（TCC 提示）
  \item \textbf{CLI}（可选）
  \item \textbf{新手引导聊天}（专用会话）
  \item 就绪
\end{enumerate}


\subsection{1) 欢迎 + 安全提示}


阅读显示的安全提示并相应决定。

\subsection{2) 本地 vs 远程}


\textbf{Gateway 网关}在哪里运行？

\begin{itemize}
  \item \textbf{本地（此 Mac）：} 新手引导可以在本地运行 OAuth 流程并写入凭证。
  \item \textbf{远程（通过 SSH/Tailnet）：} 新手引导\textbf{不会}在本地运行 OAuth；凭证必须存在于 Gateway 网关主机上。
  \item \textbf{稍后配置：} 跳过设置并保持应用未配置状态。
\end{itemize}


Gateway 网关认证提示：

\begin{itemize}
  \item 向导现在即使对于 loopback 也会生成\textbf{令牌}，因此本地 WS 客户端必须认证。
  \item 如果你禁用认证，任何本地进程都可以连接；仅在完全受信任的机器上使用。
  \item 对于多机器访问或非 loopback 绑定，使用\textbf{令牌}。
\end{itemize}


\subsection{3) 仅限本地的认证（Anthropic OAuth）}


macOS 应用支持 Anthropic OAuth（Claude Pro/Max）。流程：

\begin{itemize}
  \item 打开浏览器进行 OAuth（PKCE）
  \item 要求用户粘贴 \texttt{code\#state} 值
  \item 将凭证写入 \texttt{\textasciitilde{}/.openclaw/credentials/oauth.json}
\end{itemize}


其他提供商（OpenAI、自定义 API）目前通过环境变量或配置文件配置。

\subsection{4) 设置向导（Gateway 网关驱动）}


应用可以运行与 CLI 相同的设置向导。这使新手引导与 Gateway 网关端行为保持同步，避免在 SwiftUI 中重复逻辑。

\subsection{5) 权限}


新手引导请求以下所需的 TCC 权限：

\begin{itemize}
  \item 通知
  \item 辅助功能
  \item 屏幕录制
  \item 麦克风 / 语音识别
  \item 自动化（AppleScript）
\end{itemize}


\subsection{6) CLI（可选）}


应用可以通过 npm/pnpm 安装全局 \texttt{openclaw} CLI，以便终端工作流和 launchd 任务开箱即用。

\subsection{7) 新手引导聊天（专用会话）}


设置完成后，应用会打开一个专用的新手引导聊天会话，让智能体可以自我介绍并指导后续步骤。这使首次运行指导与你的正常对话分开。

\subsection{智能体引导仪式}


在首次智能体运行时，OpenClaw 会引导一个工作区（默认 \texttt{\textasciitilde{}/.openclaw/workspace}）：

\begin{itemize}
  \item 初始化 \texttt{AGENTS.md}、\texttt{BOOTSTRAP.md}、\texttt{IDENTITY.md}、\texttt{USER.md}
  \item 运行简短的问答仪式（一次一个问题）
  \item 将身份 + 偏好写入 \texttt{IDENTITY.md}、\texttt{USER.md}、\texttt{SOUL.md}
  \item 完成后删除 \texttt{BOOTSTRAP.md}，使其只运行一次
\end{itemize}


\subsection{可选：Gmail 钩子（手动）}


Gmail Pub/Sub 设置目前是手动步骤。使用：

\begin{lstlisting}[language=bash]
openclaw webhooks gmail setup --account you@gmail.com
\end{lstlisting}


参阅 \href{/automation/gmail-pubsub}{/automation/gmail-pubsub} 了解详情。

\subsection{远程模式说明}


当 Gateway 网关在另一台机器上运行时，凭证和工作区文件存储在\textbf{该主机上}。如果你需要在远程模式下使用 OAuth，请在 Gateway 网关主机上创建：

\begin{itemize}
  \item \texttt{\textasciitilde{}/.openclaw/credentials/oauth.json}
  \item \texttt{\textasciitilde{}/.openclaw/agents//agent/auth-profiles.json}
\end{itemize}



\clearpage

% === 源文件: start/openclaw.md ===
\section{使用 OpenClaw 构建个人助手}


OpenClaw 是 \textbf{Pi} 智能体的 WhatsApp + Telegram + Discord + iMessage Gateway 网关。插件可添加 Mattermost。本指南是"个人助手"设置：一个专用的 WhatsApp 号码，表现得像你的常驻智能体。

\subsection{\emoji{⚠} 安全第一}


你正在让智能体处于可以：

\begin{itemize}
  \item 在你的机器上运行命令（取决于你的 Pi 工具设置）
  \item 在你的工作区读/写文件
  \item 通过 WhatsApp/Telegram/Discord/Mattermost（插件）发送消息
\end{itemize}


从保守开始：

\begin{itemize}
  \item 始终设置 \texttt{channels.whatsapp.allowFrom}（永远不要在你的个人 Mac 上对全世界开放）。
  \item 为助手使用专用的 WhatsApp 号码。
  \item 心跳现在默认每 30 分钟一次。在你信任设置之前，通过设置 \texttt{agents.defaults.heartbeat.every: "0m"} 来禁用。
\end{itemize}


\subsection{先决条件}


\begin{itemize}
  \item Node \textbf{22+}
  \item OpenClaw 在 PATH 中可用（推荐：全局安装）
  \item 助手的第二个手机号码（SIM/eSIM/预付费）
\end{itemize}


\begin{lstlisting}[language=bash]
npm install -g openclaw@latest
# 或：pnpm add -g openclaw@latest
\end{lstlisting}


从源代码（开发）：

\begin{lstlisting}[language=bash]
git clone https://github.com/openclaw/openclaw.git
cd openclaw
pnpm install
pnpm ui:build # 首次运行时自动安装 UI 依赖
pnpm build
pnpm link --global
\end{lstlisting}


\subsection{双手机设置（推荐）}


你需要这样：

\begin{lstlisting}
你的手机（个人）               第二部手机（助手）
┌─────────────────┐           ┌─────────────────┐
│  你的 WhatsApp  │  ──────▶  │   助手 WA       │
│  +1-555-YOU     │  消息     │  +1-555-ASSIST  │
└─────────────────┘           └────────┬────────┘
                                       │ 通过二维码关联
                                       ▼
                              ┌─────────────────┐
                              │  你的 Mac       │
                              │  (openclaw)     │
                              │    Pi 智能体    │
                              └─────────────────┘
\end{lstlisting}


如果你将个人 WhatsApp 关联到 OpenClaw，发给你的每条消息都会变成"智能体输入"。这通常不是你想要的。

\subsection{5 分钟快速开始}


\begin{enumerate}
  \item 配对 WhatsApp Web（显示二维码；用助手手机扫描）：
\end{enumerate}


\begin{lstlisting}[language=bash]
openclaw channels login
\end{lstlisting}


\begin{enumerate}
  \item 启动 Gateway 网关（保持运行）：
\end{enumerate}


\begin{lstlisting}[language=bash]
openclaw gateway --port 18789
\end{lstlisting}


\begin{enumerate}
  \item 在 \texttt{\textasciitilde{}/.openclaw/openclaw.json} 中放置最小配置：
\end{enumerate}


\begin{lstlisting}[language=json]
{
  channels: { whatsapp: { allowFrom: ["+15555550123"] } },
}
\end{lstlisting}


现在从你允许列表中的手机向助手号码发消息。

新手引导完成后，我们会自动打开带有 Gateway 网关令牌的仪表板并打印带令牌的链接。稍后重新打开：\texttt{openclaw dashboard}。

\subsection{给智能体一个工作区（AGENTS）}


OpenClaw 从其工作区目录读取操作指令和"记忆"。

默认情况下，OpenClaw 使用 \texttt{\textasciitilde{}/.openclaw/workspace} 作为智能体工作区，并会在设置/首次智能体运行时自动创建它（加上起始的 \texttt{AGENTS.md}、\texttt{SOUL.md}、\texttt{TOOLS.md}、\texttt{IDENTITY.md}、\texttt{USER.md}）。\texttt{BOOTSTRAP.md} 仅在工作区是全新的时候创建（删除后不应再出现）。

提示：将此文件夹视为 OpenClaw 的"记忆"，并将其设为 git 仓库（最好是私有的），这样你的 \texttt{AGENTS.md} + 记忆文件就有了备份。如果安装了 git，全新的工作区会自动初始化。

\begin{lstlisting}[language=bash]
openclaw setup
\end{lstlisting}


完整工作区布局 + 备份指南：\href{/concepts/agent-workspace}{智能体工作区}
记忆工作流：\href{/concepts/memory}{记忆}

可选：使用 \texttt{agents.defaults.workspace} 选择不同的工作区（支持 \texttt{\textasciitilde{}}）。

\begin{lstlisting}[language=json]
{
  agent: {
    workspace: "~/.openclaw/workspace",
  },
}
\end{lstlisting}


如果你已经从仓库提供了自己的工作区文件，可以完全禁用引导文件创建：

\begin{lstlisting}[language=json]
{
  agent: {
    skipBootstrap: true,
  },
}
\end{lstlisting}


\subsection{将其变成"助手"的配置}


OpenClaw 默认为良好的助手设置，但你通常需要调整：

\begin{itemize}
  \item \texttt{SOUL.md} 中的人设/指令
  \item 思考默认值（如果需要）
  \item 心跳（一旦你信任它）
\end{itemize}


示例：

\begin{lstlisting}[language=json]
{
  logging: { level: "info" },
  agent: {
    model: "anthropic/claude-opus-4-5",
    workspace: "~/.openclaw/workspace",
    thinkingDefault: "high",
    timeoutSeconds: 1800,
    // 从 0 开始；稍后启用。
    heartbeat: { every: "0m" },
  },
  channels: {
    whatsapp: {
      allowFrom: ["+15555550123"],
      groups: {
        "*": { requireMention: true },
      },
    },
  },
  routing: {
    groupChat: {
      mentionPatterns: ["@openclaw", "openclaw"],
    },
  },
  session: {
    scope: "per-sender",
    resetTriggers: ["/new", "/reset"],
    reset: {
      mode: "daily",
      atHour: 4,
      idleMinutes: 10080,
    },
  },
}
\end{lstlisting}


\subsection{会话和记忆}


\begin{itemize}
  \item 会话文件：\texttt{\textasciitilde{}/.openclaw/agents//sessions/\{\{SessionId\}\}.jsonl}
  \item 会话元数据（token 使用量、最后路由等）：\texttt{\textasciitilde{}/.openclaw/agents//sessions/sessions.json}（旧版：\texttt{\textasciitilde{}/.openclaw/sessions/sessions.json}）
  \item \texttt{/new} 或 \texttt{/reset} 为该聊天启动新会话（可通过 \texttt{resetTriggers} 配置）。如果单独发送，智能体会回复一个简短的问候来确认重置。
  \item \texttt{/compact [instructions]} 压缩会话上下文并报告剩余的上下文预算。
\end{itemize}


\subsection{心跳（主动模式）}


默认情况下，OpenClaw 每 30 分钟运行一次心跳，提示词为：
\texttt{Read HEARTBEAT.md if it exists (workspace context). Follow it strictly. Do not infer or repeat old tasks from prior chats. If nothing needs attention, reply HEARTBEAT\_OK.}
设置 \texttt{agents.defaults.heartbeat.every: "0m"} 来禁用。

\begin{itemize}
  \item 如果 \texttt{HEARTBEAT.md} 存在但实际上是空的（只有空行和 markdown 标题如 \texttt{\# Heading}），OpenClaw 会跳过心跳运行以节省 API 调用。
  \item 如果文件不存在，心跳仍然运行，模型决定做什么。
  \item 如果智能体回复 \texttt{HEARTBEAT\_OK}（可选带有短填充；参见 \texttt{agents.defaults.heartbeat.ackMaxChars}），OpenClaw 会为该心跳抑制出站投递。
  \item 心跳运行完整的智能体轮次 — 更短的间隔会消耗更多 token。
\end{itemize}


\begin{lstlisting}[language=json]
{
  agent: {
    heartbeat: { every: "30m" },
  },
}
\end{lstlisting}


\subsection{媒体输入和输出}


入站附件（图片/音频/文档）可以通过模板暴露给你的命令：

\begin{itemize}
  \item \texttt{\{\{MediaPath\}\}}（本地临时文件路径）
  \item \texttt{\{\{MediaUrl\}\}}（伪 URL）
  \item \texttt{\{\{Transcript\}\}}（如果启用了音频转录）
\end{itemize}


来自智能体的出站附件：在单独一行包含 \texttt{MEDIA:}（无空格）。示例：

\begin{lstlisting}
这是截图。
MEDIA:https://example.com/screenshot.png
\end{lstlisting}


OpenClaw 会提取这些并将它们作为媒体与文本一起发送。

\subsection{运维检查清单}


\begin{lstlisting}[language=bash]
openclaw status          # 本地状态（凭证、会话、排队事件）
openclaw status --all    # 完整诊断（只读，可粘贴）
openclaw status --deep   # 添加 Gateway 网关健康探测（Telegram + Discord）
openclaw health --json   # Gateway 网关健康快照（WS）
\end{lstlisting}


日志位于 \texttt{/tmp/openclaw/}（默认：\texttt{openclaw-YYYY-MM-DD.log}）。

\subsection{下一步}


\begin{itemize}
  \item WebChat：\href{/web/webchat}{WebChat}
  \item Gateway 网关运维：\href{/gateway}{Gateway 网关运行手册}
  \item 定时任务 + 唤醒：\href{/automation/cron-jobs}{定时任务}
  \item macOS 菜单栏配套应用：\href{/platforms/macos}{OpenClaw macOS 应用}
  \item iOS 节点应用：\href{/platforms/ios}{iOS 应用}
  \item Android 节点应用：\href{/platforms/android}{Android 应用}
  \item Windows 状态：\href{/platforms/windows}{Windows (WSL2)}
  \item Linux 状态：\href{/platforms/linux}{Linux 应用}
  \item 安全：\href{/gateway/security}{安全}
\end{itemize}



\clearpage

\section{快速开始}
% 源文件: start/quickstart.md

% === 源文件: start/quickstart.md ===
OpenClaw 需要 Node 22 或更新版本。

\subsection{安装}


\begin{lstlisting}[language=bash]
    npm install -g openclaw@latest
\end{lstlisting}

\begin{lstlisting}[language=bash]
    pnpm add -g openclaw@latest
\end{lstlisting}


\subsection{新手引导并运行 Gateway 网关}


\begin{lstlisting}[language=bash]
    openclaw onboard --install-daemon
\end{lstlisting}

\begin{lstlisting}[language=bash]
    openclaw channels login
\end{lstlisting}

\begin{lstlisting}[language=bash]
    openclaw gateway --port 18789
\end{lstlisting}


完成新手引导后，Gateway 网关将通过用户服务运行。你也可以使用 \texttt{openclaw gateway} 手动启动。

之后在 npm 安装和 git 安装之间切换非常简单。安装另一种方式后，运行
\texttt{openclaw doctor} 即可更新 Gateway 网关服务入口点。

\subsection{从源码安装（开发）}


\begin{lstlisting}[language=bash]
git clone https://github.com/openclaw/openclaw.git
cd openclaw
pnpm install
pnpm ui:build # 首次运行时会自动安装 UI 依赖
pnpm build
openclaw onboard --install-daemon
\end{lstlisting}


如果你还没有全局安装，可以在仓库目录中通过 \texttt{pnpm openclaw ...} 运行新手引导。

\subsection{多实例快速开始（可选）}


\begin{lstlisting}[language=bash]
OPENCLAW_CONFIG_PATH=~/.openclaw/a.json \
OPENCLAW_STATE_DIR=~/.openclaw-a \
openclaw gateway --port 19001
\end{lstlisting}


\subsection{发送测试消息}


需要一个正在运行的 Gateway 网关。

\begin{lstlisting}[language=bash]
openclaw message send --target +15555550123 --message "Hello from OpenClaw"
\end{lstlisting}



\clearpage

% === 源文件: start/setup.md ===
\section{设置}


最后更新：2026-01-01

\subsection{太长不看}


\begin{itemize}
  \item \textbf{个性化设置存放在仓库之外：} \texttt{\textasciitilde{}/.openclaw/workspace}（工作区）+ \texttt{\textasciitilde{}/.openclaw/openclaw.json}（配置）。
  \item \textbf{稳定工作流：} 安装 macOS 应用；让它运行内置的 Gateway 网关。
  \item \textbf{前沿工作流：} 通过 \texttt{pnpm gateway:watch} 自己运行 Gateway 网关，然后让 macOS 应用以本地模式连接。
\end{itemize}


\subsection{先决条件（从源码）}


\begin{itemize}
  \item Node \texttt{>=22}
  \item \texttt{pnpm}
  \item Docker（可选；仅用于容器化设置/e2e — 参阅 \href{/install/docker}{Docker}）
\end{itemize}


\subsection{个性化策略（让更新不会造成问题）}


如果你想要"100\% 为我定制"\textit{并且}易于更新，将你的自定义内容保存在：

\begin{itemize}
  \item \textbf{配置：} \texttt{\textasciitilde{}/.openclaw/openclaw.json}（JSON/JSON5 格式）
  \item \textbf{工作区：} \texttt{\textasciitilde{}/.openclaw/workspace}（Skills、提示、记忆；将其设为私有 git 仓库）
\end{itemize}


引导一次：

\begin{lstlisting}[language=bash]
openclaw setup
\end{lstlisting}


在此仓库内部，使用本地 CLI 入口：

\begin{lstlisting}[language=bash]
openclaw setup
\end{lstlisting}


如果你还没有全局安装，通过 \texttt{pnpm openclaw setup} 运行它。

\subsection{稳定工作流（macOS 应用优先）}


\begin{enumerate}
  \item 安装并启动 \textbf{OpenClaw.app}（菜单栏）。
  \item 完成新手引导/权限检查清单（TCC 提示）。
  \item 确保 Gateway 网关是\textbf{本地}并正在运行（应用管理它）。
  \item 链接表面（示例：WhatsApp）：
\end{enumerate}


\begin{lstlisting}[language=bash]
openclaw channels login
\end{lstlisting}


\begin{enumerate}
  \item 完整性检查：
\end{enumerate}


\begin{lstlisting}[language=bash]
openclaw health
\end{lstlisting}


如果你的构建版本中没有新手引导：

\begin{itemize}
  \item 运行 \texttt{openclaw setup}，然后 \texttt{openclaw channels login}，然后手动启动 Gateway 网关（\texttt{openclaw gateway}）。
\end{itemize}


\subsection{前沿工作流（在终端中运行 Gateway 网关）}


目标：开发 TypeScript Gateway 网关，获得热重载，保持 macOS 应用 UI 连接。

\subsubsection{0)（可选）也从源码运行 macOS 应用}


如果你也想让 macOS 应用保持前沿：

\begin{lstlisting}[language=bash]
./scripts/restart-mac.sh
\end{lstlisting}


\subsubsection{1) 启动开发 Gateway 网关}


\begin{lstlisting}[language=bash]
pnpm install
pnpm gateway:watch
\end{lstlisting}


\texttt{gateway:watch} 以监视模式运行 Gateway 网关，并在 TypeScript 更改时重新加载。

\subsubsection{2) 将 macOS 应用指向你正在运行的 Gateway 网关}


在 \textbf{OpenClaw.app} 中：

\begin{itemize}
  \item 连接模式：\textbf{本地}
\end{itemize}

应用将连接到在配置端口上运行的 Gateway 网关。

\subsubsection{3) 验证}


\begin{itemize}
  \item 应用内 Gateway 网关状态应显示 \textbf{"Using existing gateway …"}
  \item 或通过 CLI：
\end{itemize}


\begin{lstlisting}[language=bash]
openclaw health
\end{lstlisting}


\subsubsection{常见陷阱}


\begin{itemize}
  \item \textbf{端口错误：} Gateway 网关 WS 默认为 \texttt{ws://127.0.0.1:18789}；保持应用 + CLI 在同一端口上。
  \item \textbf{状态存储位置：}
  \item 凭证：\texttt{\textasciitilde{}/.openclaw/credentials/}
  \item 会话：\texttt{\textasciitilde{}/.openclaw/agents//sessions/}
  \item 日志：\texttt{/tmp/openclaw/}
\end{itemize}


\subsection{凭证存储映射}


在调试认证或决定备份什么时使用此映射：

\begin{itemize}
  \item \textbf{WhatsApp}：\texttt{\textasciitilde{}/.openclaw/credentials/whatsapp//creds.json}
  \item \textbf{Telegram bot token}：配置/环境变量或 \texttt{channels.telegram.tokenFile}
  \item \textbf{Discord bot token}：配置/环境变量（尚不支持令牌文件）
  \item \textbf{Slack tokens}：配置/环境变量（\texttt{channels.slack.*}）
  \item \textbf{配对允许列表}：\texttt{\textasciitilde{}/.openclaw/credentials/-allowFrom.json}
  \item \textbf{模型认证配置文件}：\texttt{\textasciitilde{}/.openclaw/agents//agent/auth-profiles.json}
  \item \textbf{旧版 OAuth 导入}：\texttt{\textasciitilde{}/.openclaw/credentials/oauth.json}
\end{itemize}

更多详情：\href{/gateway/security\#credential-storage-map}{安全}。

\subsection{更新（不破坏你的设置）}


\begin{itemize}
  \item 将 \texttt{\textasciitilde{}/.openclaw/workspace} 和 \texttt{\textasciitilde{}/.openclaw/} 保持为"你的东西"；不要将个人提示/配置放入 \texttt{openclaw} 仓库。
  \item 更新源码：\texttt{git pull} + \texttt{pnpm install}（当锁文件更改时）+ 继续使用 \texttt{pnpm gateway:watch}。
\end{itemize}


\subsection{Linux（systemd 用户服务）}


Linux 安装使用 systemd \textbf{用户}服务。默认情况下，systemd 在注销/空闲时停止用户服务，这会终止 Gateway 网关。新手引导会尝试为你启用 lingering（可能提示 sudo）。如果仍然关闭，运行：

\begin{lstlisting}[language=bash]
sudo loginctl enable-linger $USER
\end{lstlisting}


对于常驻或多用户服务器，考虑使用\textbf{系统}服务而不是用户服务（不需要 lingering）。参阅 \href{/gateway}{Gateway 网关运行手册} 了解 systemd 说明。

\subsection{相关文档}


\begin{itemize}
  \item \href{/gateway}{Gateway 网关运行手册}（标志、监督、端口）
  \item \href{/gateway/configuration}{Gateway 网关配置}（配置模式 + 示例）
  \item \href{/channels/discord}{Discord} 和 \href{/channels/telegram}{Telegram}（回复标签 + replyToMode 设置）
  \item \href{/start/openclaw}{OpenClaw 助手设置}
  \item \href{/platforms/macos}{macOS 应用}（Gateway 网关生命周期）
\end{itemize}



\clearpage

% === 源文件: start/showcase.md ===
\section{案例展示}


来自社区的真实项目。看看大家正在用 OpenClaw 构建什么。

\textbf{想要展示你的项目？} 在 \href{https://discord.gg/clawd}{Discord 的 \#showcase 频道} 分享或在 \href{https://x.com/openclaw}{X 上 @openclaw}。

\subsection{OpenClaw 实战演示}


VelvetShark 的完整设置演练（28 分钟）。

<div
position: "relative",
paddingBottom: "56.25\%",
height: 0,
overflow: "hidden",
borderRadius: 16,
\}\}
\begin{quote}

\end{quote}

<iframe
title="OpenClaw: The self-hosted AI that Siri should have been (Full setup)"
frameBorder="0"
loading="lazy"
allow="accelerometer; autoplay; clipboard-write; encrypted-media; gyroscope; picture-in-picture; web-share"
allowFullScreen

\href{https://www.youtube.com/watch?v=SaWSPZoPX34}{在 YouTube 上观看}

<div
position: "relative",
paddingBottom: "56.25\%",
height: 0,
overflow: "hidden",
borderRadius: 16,
\}\}
\begin{quote}

\end{quote}

<iframe
title="OpenClaw showcase video"
frameBorder="0"
loading="lazy"
allow="accelerometer; autoplay; clipboard-write; encrypted-media; gyroscope; picture-in-picture; web-share"
allowFullScreen

\href{https://www.youtube.com/watch?v=mMSKQvlmFuQ}{在 YouTube 上观看}

<div
position: "relative",
paddingBottom: "56.25\%",
height: 0,
overflow: "hidden",
borderRadius: 16,
\}\}
\begin{quote}

\end{quote}

<iframe
title="OpenClaw community showcase"
frameBorder="0"
loading="lazy"
allow="accelerometer; autoplay; clipboard-write; encrypted-media; gyroscope; picture-in-picture; web-share"
allowFullScreen

\href{https://www.youtube.com/watch?v=5kkIJNUGFho}{在 YouTube 上观看}

\subsection{\emoji{🆕} Discord 最新分享}



\textbf{@bangnokia} • \texttt{review} \texttt{github} \texttt{telegram}

OpenCode 完成更改 → 打开 PR → OpenClaw 审查差异并在 Telegram 中回复"小建议"加上明确的合并决定（包括需要先应用的关键修复）。


\textbf{@prades\_maxime} • \texttt{skills} \texttt{local} \texttt{csv}

向"Robby"（@openclaw）请求一个本地酒窖 skill。它请求一个示例 CSV 导出 + 存储位置，然后快速构建/测试该 skill（示例中有 962 瓶酒）。


\textbf{@marchattonhere} • \texttt{automation} \texttt{browser} \texttt{shopping}

每周餐饮计划 → 常购商品 → 预订配送时段 → 确认订单。无需 API，仅使用浏览器控制。


\textbf{@am-will} • \texttt{devtools} \texttt{screenshots} \texttt{markdown}

快捷键选择屏幕区域 → Gemini 视觉 → 即时 Markdown 到剪贴板。


\textbf{@kitze} • \texttt{ui} \texttt{skills} \texttt{sync}

跨 Agents、Claude、Codex 和 OpenClaw 管理 skills/命令的桌面应用。


\textbf{社区} • \texttt{voice} \texttt{tts} \texttt{telegram}

封装 papla.media TTS 并将结果作为 Telegram 语音备忘录发送（无烦人的自动播放）。


\textbf{@odrobnik} • \texttt{devtools} \texttt{codex} \texttt{brew}

Homebrew 安装的助手工具，用于列出/检查/监视本地 OpenAI Codex 会话（CLI + VS Code）。


\textbf{@tobiasbischoff} • \texttt{hardware} \texttt{3d-printing} \texttt{skill}

控制和排查 BambuLab 打印机：状态、任务、摄像头、AMS、校准等。


\textbf{@hjanuschka} • \texttt{travel} \texttt{transport} \texttt{skill}

维也纳公共交通的实时发车时间、中断信息、电梯状态和路线规划。


\textbf{@George5562} • \texttt{automation} \texttt{browser} \texttt{parenting}

通过 ParentPay 自动预订英国学校餐食。使用鼠标坐标实现可靠的表格单元格点击。

\textbf{@julianengel} • \texttt{files} \texttt{r2} \texttt{presigned-urls}

上传到 Cloudflare R2/S3 并生成安全的预签名下载链接。非常适合远程 OpenClaw 实例。

\textbf{@coard} • \texttt{ios} \texttt{xcode} \texttt{testflight}

构建了一个完整的带有地图和语音录制功能的 iOS 应用，完全通过 Telegram 聊天部署到 TestFlight。


\textbf{@AS} • \texttt{health} \texttt{oura} \texttt{calendar}

个人 AI 健康助手，将 Oura 戒指数据与日历、预约和健身房计划集成。

\textbf{@adam91holt} • \texttt{multi-agent} \texttt{orchestration} \texttt{architecture} \texttt{manifesto}

一个 Gateway 网关下的 14+ 智能体，Opus 4.5 编排器将任务委派给 Codex 工作者。全面的\href{https://github.com/adam91holt/orchestrated-ai-articles}{技术文章}涵盖梦之队阵容、模型选择、沙箱隔离、webhook、心跳和委派流程。用于智能体沙箱隔离的 \href{https://github.com/adam91holt/clawdspace}{Clawdspace}。\href{https://adams-ai-journey.ghost.io/2026-the-year-of-the-orchestrator/}{博客文章}。

\textbf{@NessZerra} • \texttt{devtools} \texttt{linear} \texttt{cli} \texttt{issues}

与智能体工作流（Claude Code、OpenClaw）集成的 Linear CLI。从终端管理问题、项目和工作流。首个外部 PR 已合并！

\textbf{@jules} • \texttt{messaging} \texttt{beeper} \texttt{cli} \texttt{automation}

通过 Beeper Desktop 读取、发送和归档消息。使用 Beeper 本地 MCP API，让智能体可以在一个地方管理你的所有聊天（iMessage、WhatsApp 等）。


\subsection{自动化与工作流}



\textbf{@antonplex} • \texttt{automation} \texttt{hardware} \texttt{air-quality}

Claude Code 发现并确认了净化器控制，然后 OpenClaw 接管来管理房间空气质量。


\textbf{@signalgaining} • \texttt{automation} \texttt{camera} \texttt{skill} \texttt{images}

由屋顶摄像头触发：让 OpenClaw 在天空看起来很美的时候拍一张照片——它设计了一个 skill 并拍摄了照片。


\textbf{@buddyhadry} • \texttt{automation} \texttt{briefing} \texttt{images} \texttt{telegram}

定时提示每天早上通过 OpenClaw 角色生成一张"场景"图片（天气、任务、日期、喜欢的帖子/引言）。

\textbf{@joshp123} • \texttt{automation} \texttt{booking} \texttt{cli}

Playtomic 可用性检查器 + 预订 CLI。再也不会错过空闲场地。


\textbf{社区} • \texttt{automation} \texttt{email} \texttt{pdf}

从邮件收集 PDF，为税务顾问准备文档。月度会计自动运行。

\textbf{@davekiss} • \texttt{telegram} \texttt{website} \texttt{migration} \texttt{astro}

一边看 Netflix 一边通过 Telegram 重建整个个人网站——Notion → Astro，迁移了 18 篇文章，DNS 转到 Cloudflare。从未打开笔记本电脑。

\textbf{@attol8} • \texttt{automation} \texttt{api} \texttt{skill}

搜索职位列表，与简历关键词匹配，返回带链接的相关机会。使用 JSearch API 在 30 分钟内构建。

\textbf{@jdrhyne} • \texttt{automation} \texttt{jira} \texttt{skill} \texttt{devtools}

OpenClaw 连接到 Jira，然后即时生成一个新的 skill（在它出现在 ClawHub 之前）。

\textbf{@iamsubhrajyoti} • \texttt{automation} \texttt{todoist} \texttt{skill} \texttt{telegram}

自动化 Todoist 任务，并让 OpenClaw 直接在 Telegram 聊天中生成 skill。

\textbf{@bheem1798} • \texttt{finance} \texttt{browser} \texttt{automation}

通过浏览器自动化登录 TradingView，截取图表屏幕截图，并按需执行技术分析。无需 API——只需浏览器控制。

\textbf{@henrymascot} • \texttt{slack} \texttt{automation} \texttt{support}

监视公司 Slack 频道，提供有用的回复，并将通知转发到 Telegram。在没有被要求的情况下自主修复了已部署应用中的生产 bug。


\subsection{知识与记忆}



\textbf{@joshp123} • \texttt{learning} \texttt{voice} \texttt{skill}

通过 OpenClaw 实现带有发音反馈和学习流程的中文学习引擎。


\textbf{社区} • \texttt{memory} \texttt{transcription} \texttt{indexing}

导入完整的 WhatsApp 导出，转录 1k+ 条语音备忘录，与 git 日志交叉检查，输出链接的 markdown 报告。

\textbf{@jamesbrooksco} • \texttt{search} \texttt{vector} \texttt{bookmarks}

使用 Qdrant + OpenAI/Ollama embeddings 为 Karakeep 书签添加向量搜索。

\textbf{社区} • \texttt{memory} \texttt{beliefs} \texttt{self-model}

独立的记忆管理器，将会话文件转化为记忆 → 信念 → 演化的自我模型。


\subsection{语音与电话}



\textbf{@alejandroOPI} • \texttt{voice} \texttt{vapi} \texttt{bridge}

Vapi 语音助手 $\leftrightarrow$ OpenClaw HTTP 桥接。与你的智能体进行近实时电话通话。

\textbf{@obviyus} • \texttt{transcription} \texttt{multilingual} \texttt{skill}

通过 OpenRouter（Gemini 等）进行多语言音频转录。可在 ClawHub 获取。


\subsection{基础设施与部署}



\textbf{@ngutman} • \texttt{homeassistant} \texttt{docker} \texttt{raspberry-pi}

在 Home Assistant OS 上运行的 OpenClaw Gateway 网关，支持 SSH 隧道和持久状态。

\textbf{ClawHub} • \texttt{homeassistant} \texttt{skill} \texttt{automation}

通过自然语言控制和自动化 Home Assistant 设备。

\textbf{@openclaw} • \texttt{nix} \texttt{packaging} \texttt{deployment}

开箱即用的 nixified OpenClaw 配置，用于可复现的部署。

\textbf{ClawHub} • \texttt{calendar} \texttt{caldav} \texttt{skill}

使用 khal/vdirsyncer 的日历 skill。自托管日历集成。


\subsection{家居与硬件}



\textbf{@joshp123} • \texttt{home} \texttt{nix} \texttt{grafana}

Nix 原生家庭自动化，以 OpenClaw 作为界面，加上漂亮的 Grafana 仪表板。


\textbf{@joshp123} • \texttt{vacuum} \texttt{iot} \texttt{plugin}

通过自然对话控制你的 Roborock 扫地机器人。



\subsection{社区项目}



\textbf{社区} • \texttt{marketplace} \texttt{astronomy} \texttt{webapp}

完整的天文设备市场。围绕 OpenClaw 生态系统构建。


\hrulefill


\subsection{提交你的项目}


有想分享的东西？我们很乐意展示它！

在 \href{https://discord.gg/clawd}{Discord 的 \#showcase 频道} 发布或在 \href{https://x.com/openclaw}{Twitter 上 @openclaw}
告诉我们它做什么，链接到仓库/演示，如果有的话分享截图
我们会将优秀项目添加到此页面


\clearpage

% === 源文件: start/wizard.md ===
\section{新手引导向导（CLI）}


新手引导向导是在 macOS、Linux 或 Windows（通过 WSL2；强烈推荐）上设置 OpenClaw 的\textbf{推荐}方式。
它可以在一个引导式流程中配置本地 Gateway 网关或远程 Gateway 网关连接，以及渠道、Skills 和工作区默认值。

主要入口：

\begin{lstlisting}[language=bash]
openclaw onboard
\end{lstlisting}


最快开始聊天的方式：打开控制界面（无需设置渠道）。运行 \texttt{openclaw dashboard} 并在浏览器中聊天。文档：\href{/web/dashboard}{控制面板}。

后续重新配置：

\begin{lstlisting}[language=bash]
openclaw configure
\end{lstlisting}


推荐：设置 Brave Search API 密钥，以便智能体可以使用 \texttt{web\_search}（\texttt{web\_fetch} 无需密钥即可使用）。最简单的方式：\texttt{openclaw configure --section web}，它会存储 \texttt{tools.web.search.apiKey}。文档：\href{/tools/web}{Web 工具}。

\subsection{快速开始 vs 高级}


向导从\textbf{快速开始}（默认值）vs \textbf{高级}（完全控制）开始。

\textbf{快速开始}保持默认值：

\begin{itemize}
  \item 本地 Gateway 网关（loopback）
  \item 默认工作区（或现有工作区）
  \item Gateway 网关端口 \textbf{18789}
  \item Gateway 网关认证 \textbf{Token}（自动生成，即使在 loopback 上）
  \item Tailscale 暴露 \textbf{关闭}
  \item Telegram + WhatsApp 私信默认使用\textbf{允许列表}（系统会提示你输入电话号码）
\end{itemize}


\textbf{高级}暴露每个步骤（模式、工作区、Gateway 网关、渠道、守护进程、Skills）。

\subsection{向导做了什么}


\textbf{本地模式（默认）}引导你完成：

\begin{itemize}
  \item 模型/认证（OpenAI Code (Codex) 订阅 OAuth、Anthropic API 密钥（推荐）或 setup-token（粘贴），以及 MiniMax/GLM/Moonshot/AI Gateway 选项）
  \item 工作区位置 + 引导文件
  \item Gateway 网关设置（端口/绑定/认证/tailscale）
  \item 提供商（Telegram、WhatsApp、Discord、Google Chat、Mattermost（插件）、Signal）
  \item 守护进程安装（LaunchAgent / systemd 用户单元）
  \item 健康检查
  \item Skills（推荐）
\end{itemize}


\textbf{远程模式}仅配置本地客户端连接到其他位置的 Gateway 网关。
它\textbf{不会}在远程主机上安装或更改任何内容。

要添加更多隔离的智能体（独立的工作区 + 会话 + 认证），使用：

\begin{lstlisting}[language=bash]
openclaw agents add <name>
\end{lstlisting}


提示：\texttt{--json} \textbf{不}意味着非交互模式。脚本中请使用 \texttt{--non-interactive}（和 \texttt{--workspace}）。

\subsection{流程详情（本地）}


\begin{enumerate}
  \item \textbf{现有配置检测}
\end{enumerate}

\begin{itemize}
  \item 如果 \texttt{\textasciitilde{}/.openclaw/openclaw.json} 存在，选择\textbf{保留 / 修改 / 重置}。
  \item 重新运行向导\textbf{不会}清除任何内容，除非你明确选择\textbf{重置}（或传递 \texttt{--reset}）。
  \item 如果配置无效或包含遗留键名，向导会停止并要求你在继续之前运行 \texttt{openclaw doctor}。
  \item 重置使用 \texttt{trash}（永不使用 \texttt{rm}）并提供范围选项：
  \item 仅配置
  \item 配置 + 凭证 + 会话
  \item 完全重置（同时删除工作区）
\end{itemize}


\begin{enumerate}
  \item \textbf{模型/认证}
\end{enumerate}

\begin{itemize}
  \item \textbf{Anthropic API 密钥（推荐）}：如果存在则使用 \texttt{ANTHROPIC\_API\_KEY}，否则提示输入密钥，然后保存供守护进程使用。
  \item \textbf{Anthropic OAuth（Claude Code CLI）}：在 macOS 上，向导检查钥匙串项目"Claude Code-credentials"（选择"始终允许"以便 launchd 启动不会阻塞）；在 Linux/Windows 上，如果存在则复用 \texttt{\textasciitilde{}/.claude/.credentials.json}。
  \item \textbf{Anthropic 令牌（粘贴 setup-token）}：在任何机器上运行 \texttt{claude setup-token}，然后粘贴令牌（你可以命名它；空白 = 默认）。
  \item \textbf{OpenAI Code (Codex) 订阅（Codex CLI）}：如果 \texttt{\textasciitilde{}/.codex/auth.json} 存在，向导可以复用它。
  \item \textbf{OpenAI Code (Codex) 订阅（OAuth）}：浏览器流程；粘贴 \texttt{code\#state}。
  \item 当模型未设置或为 \texttt{openai/*} 时，将 \texttt{agents.defaults.model} 设置为 \texttt{openai-codex/gpt-5.2}。
  \item \textbf{OpenAI API 密钥}：如果存在则使用 \texttt{OPENAI\_API\_KEY}，否则提示输入密钥，然后保存到 \texttt{\textasciitilde{}/.openclaw/.env} 以便 launchd 可以读取。
  \item \textbf{OpenCode Zen（多模型代理）}：提示输入 \texttt{OPENCODE\_API\_KEY}（或 \texttt{OPENCODE\_ZEN\_API\_KEY}，在 https://opencode.ai/auth 获取）。
  \item \textbf{API 密钥}：为你存储密钥。
  \item \textbf{Vercel AI Gateway（多模型代理）}：提示输入 \texttt{AI\_GATEWAY\_API\_KEY}。
  \item 更多详情：\href{/providers/vercel-ai-gateway}{Vercel AI Gateway}
  \item \textbf{MiniMax M2.1}：自动写入配置。
  \item 更多详情：\href{/providers/minimax}{MiniMax}
  \item \textbf{Synthetic（Anthropic 兼容）}：提示输入 \texttt{SYNTHETIC\_API\_KEY}。
  \item 更多详情：\href{/providers/synthetic}{Synthetic}
  \item \textbf{Moonshot（Kimi K2）}：自动写入配置。
  \item \textbf{Kimi Coding}：自动写入配置。
  \item 更多详情：\href{/providers/moonshot}{Moonshot AI（Kimi + Kimi Coding）}
  \item \textbf{跳过}：尚未配置认证。
  \item 从检测到的选项中选择默认模型（或手动输入提供商/模型）。
  \item 向导运行模型检查，如果配置的模型未知或缺少认证则发出警告。
\end{itemize}


\begin{itemize}
  \item OAuth 凭证存储在 \texttt{\textasciitilde{}/.openclaw/credentials/oauth.json}；认证配置文件存储在 \texttt{\textasciitilde{}/.openclaw/agents//agent/auth-profiles.json}（API 密钥 + OAuth）。
  \item 更多详情：\href{/concepts/oauth}{/concepts/oauth}
\end{itemize}


\begin{enumerate}
  \item \textbf{工作区}
\end{enumerate}

\begin{itemize}
  \item 默认 \texttt{\textasciitilde{}/.openclaw/workspace}（可配置）。
  \item 为智能体引导仪式播种所需的工作区文件。
  \item 完整的工作区布局 + 备份指南：\href{/concepts/agent-workspace}{智能体工作区}
\end{itemize}


\begin{enumerate}
  \item \textbf{Gateway 网关}
\end{enumerate}

\begin{itemize}
  \item 端口、绑定、认证模式、tailscale 暴露。
  \item 认证建议：即使对于 loopback 也保持 \textbf{Token}，以便本地 WS 客户端必须进行认证。
  \item 仅当你完全信任每个本地进程时才禁用认证。
  \item 非 loopback 绑定仍需要认证。
\end{itemize}


\begin{enumerate}
  \item \textbf{渠道}
\end{enumerate}

\begin{itemize}
  \item \href{/channels/whatsapp}{WhatsApp}：可选的二维码登录。
  \item \href{/channels/telegram}{Telegram}：机器人令牌。
  \item \href{/channels/discord}{Discord}：机器人令牌。
  \item \href{/channels/googlechat}{Google Chat}：服务账户 JSON + webhook 受众。
  \item \href{/channels/mattermost}{Mattermost}（插件）：机器人令牌 + 基础 URL。
  \item \href{/channels/signal}{Signal}：可选的 \texttt{signal-cli} 安装 + 账户配置。
  \item \href{/channels/imessage}{iMessage}：本地 \texttt{imsg} CLI 路径 + 数据库访问。
  \item 私信安全：默认为配对。第一条私信发送验证码；通过 \texttt{openclaw pairing approve  } 批准或使用允许列表。
\end{itemize}


\begin{enumerate}
  \item \textbf{守护进程安装}
\end{enumerate}

\begin{itemize}
  \item macOS：LaunchAgent
  \item 需要已登录的用户会话；对于无头环境，使用自定义 LaunchDaemon（未提供）。
  \item Linux（和通过 WSL2 的 Windows）：systemd 用户单元
  \item 向导尝试通过 \texttt{loginctl enable-linger } 启用 lingering，以便 Gateway 网关在注销后保持运行。
  \item 可能提示 sudo（写入 \texttt{/var/lib/systemd/linger}）；它首先尝试不使用 sudo。
  \item \textbf{运行时选择：}Node（推荐；WhatsApp/Telegram 需要）。\textbf{不推荐} Bun。
\end{itemize}


\begin{enumerate}
  \item \textbf{健康检查}
\end{enumerate}

\begin{itemize}
  \item 启动 Gateway 网关（如果需要）并运行 \texttt{openclaw health}。
  \item 提示：\texttt{openclaw status --deep} 在状态输出中添加 Gateway 网关健康探测（需要可达的 Gateway 网关）。
\end{itemize}


\begin{enumerate}
  \item \textbf{Skills（推荐）}
\end{enumerate}

\begin{itemize}
  \item 读取可用的 Skills 并检查要求。
  \item 让你选择节点管理器：\textbf{npm / pnpm}（不推荐 bun）。
  \item 安装可选依赖项（某些在 macOS 上使用 Homebrew）。
\end{itemize}


\begin{enumerate}
  \item \textbf{完成}
\end{enumerate}

\begin{itemize}
  \item 总结 + 后续步骤，包括用于额外功能的 iOS/Android/macOS 应用。
\end{itemize}


\begin{itemize}
  \item 如果未检测到 GUI，向导会打印控制界面的 SSH 端口转发说明，而不是打开浏览器。
  \item 如果控制界面资源缺失，向导会尝试构建它们；回退方案是 \texttt{pnpm ui:build}（自动安装 UI 依赖）。
\end{itemize}


\subsection{远程模式}


远程模式配置本地客户端连接到其他位置的 Gateway 网关。

你将设置的内容：

\begin{itemize}
  \item 远程 Gateway 网关 URL（\texttt{ws://...}）
  \item 如果远程 Gateway 网关需要认证则需要令牌（推荐）
\end{itemize}


注意事项：

\begin{itemize}
  \item 不执行远程安装或守护进程更改。
  \item 如果 Gateway 网关仅限 loopback，使用 SSH 隧道或 tailnet。
  \item 发现提示：
  \item macOS：Bonjour（\texttt{dns-sd}）
  \item Linux：Avahi（\texttt{avahi-browse}）
\end{itemize}


\subsection{添加另一个智能体}


使用 \texttt{openclaw agents add } 创建一个具有独立工作区、会话和认证配置文件的单独智能体。不带 \texttt{--workspace} 运行会启动向导。

它设置的内容：

\begin{itemize}
  \item \texttt{agents.list[].name}
  \item \texttt{agents.list[].workspace}
  \item \texttt{agents.list[].agentDir}
\end{itemize}


注意事项：

\begin{itemize}
  \item 默认工作区遵循 \texttt{\textasciitilde{}/.openclaw/workspace-}。
  \item 添加 \texttt{bindings} 以路由入站消息（向导可以执行此操作）。
  \item 非交互标志：\texttt{--model}、\texttt{--agent-dir}、\texttt{--bind}、\texttt{--non-interactive}。
\end{itemize}


\subsection{非交互模式}


使用 \texttt{--non-interactive} 自动化或脚本化新手引导：

\begin{lstlisting}[language=bash]
openclaw onboard --non-interactive \
  --mode local \
  --auth-choice apiKey \
  --anthropic-api-key "$ANTHROPIC_API_KEY" \
  --gateway-port 18789 \
  --gateway-bind loopback \
  --install-daemon \
  --daemon-runtime node \
  --skip-skills
\end{lstlisting}


添加 \texttt{--json} 以获取机器可读的摘要。

Gemini 示例：

\begin{lstlisting}[language=bash]
openclaw onboard --non-interactive \
  --mode local \
  --auth-choice gemini-api-key \
  --gemini-api-key "$GEMINI_API_KEY" \
  --gateway-port 18789 \
  --gateway-bind loopback
\end{lstlisting}


Z.AI 示例：

\begin{lstlisting}[language=bash]
openclaw onboard --non-interactive \
  --mode local \
  --auth-choice zai-api-key \
  --zai-api-key "$ZAI_API_KEY" \
  --gateway-port 18789 \
  --gateway-bind loopback
\end{lstlisting}


Vercel AI Gateway 示例：

\begin{lstlisting}[language=bash]
openclaw onboard --non-interactive \
  --mode local \
  --auth-choice ai-gateway-api-key \
  --ai-gateway-api-key "$AI_GATEWAY_API_KEY" \
  --gateway-port 18789 \
  --gateway-bind loopback
\end{lstlisting}


Moonshot 示例：

\begin{lstlisting}[language=bash]
openclaw onboard --non-interactive \
  --mode local \
  --auth-choice moonshot-api-key \
  --moonshot-api-key "$MOONSHOT_API_KEY" \
  --gateway-port 18789 \
  --gateway-bind loopback
\end{lstlisting}


Synthetic 示例：

\begin{lstlisting}[language=bash]
openclaw onboard --non-interactive \
  --mode local \
  --auth-choice synthetic-api-key \
  --synthetic-api-key "$SYNTHETIC_API_KEY" \
  --gateway-port 18789 \
  --gateway-bind loopback
\end{lstlisting}


OpenCode Zen 示例：

\begin{lstlisting}[language=bash]
openclaw onboard --non-interactive \
  --mode local \
  --auth-choice opencode-zen \
  --opencode-zen-api-key "$OPENCODE_API_KEY" \
  --gateway-port 18789 \
  --gateway-bind loopback
\end{lstlisting}


添加智能体（非交互）示例：

\begin{lstlisting}[language=bash]
openclaw agents add work \
  --workspace ~/.openclaw/workspace-work \
  --model openai/gpt-5.2 \
  --bind whatsapp:biz \
  --non-interactive \
  --json
\end{lstlisting}


\subsection{Gateway 网关向导 RPC}


Gateway 网关通过 RPC 暴露向导流程（\texttt{wizard.start}、\texttt{wizard.next}、\texttt{wizard.cancel}、\texttt{wizard.status}）。
客户端（macOS 应用、控制界面）可以渲染步骤而无需重新实现新手引导逻辑。

\subsection{Signal 设置（signal-cli）}


向导可以从 GitHub releases 安装 \texttt{signal-cli}：

\begin{itemize}
  \item 下载适当的发布资源。
  \item 存储在 \texttt{\textasciitilde{}/.openclaw/tools/signal-cli//} 下。
  \item 将 \texttt{channels.signal.cliPath} 写入你的配置。
\end{itemize}


注意事项：

\begin{itemize}
  \item JVM 构建需要 \textbf{Java 21}。
  \item 可用时使用原生构建。
  \item Windows 使用 WSL2；signal-cli 安装在 WSL 内遵循 Linux 流程。
\end{itemize}


\subsection{向导写入的内容}


\texttt{\textasciitilde{}/.openclaw/openclaw.json} 中的典型字段：

\begin{itemize}
  \item \texttt{agents.defaults.workspace}
  \item \texttt{agents.defaults.model} / \texttt{models.providers}（如果选择了 Minimax）
  \item \texttt{gateway.*}（模式、绑定、认证、tailscale）
  \item \texttt{channels.telegram.botToken}、\texttt{channels.discord.token}、\texttt{channels.signal.\textit{}、\texttt{channels.imessage.}}
  \item 当你在提示中选择加入时的渠道允许列表（Slack/Discord/Matrix/Microsoft Teams）（名称在可能时解析为 ID）。
  \item \texttt{skills.install.nodeManager}
  \item \texttt{wizard.lastRunAt}
  \item \texttt{wizard.lastRunVersion}
  \item \texttt{wizard.lastRunCommit}
  \item \texttt{wizard.lastRunCommand}
  \item \texttt{wizard.lastRunMode}
\end{itemize}


\texttt{openclaw agents add} 写入 \texttt{agents.list[]} 和可选的 \texttt{bindings}。

WhatsApp 凭证存储在 \texttt{\textasciitilde{}/.openclaw/credentials/whatsapp//} 下。
会话存储在 \texttt{\textasciitilde{}/.openclaw/agents//sessions/} 下。

某些渠道以插件形式提供。当你在新手引导期间选择一个时，向导会在配置之前提示安装它（npm 或本地路径）。

\subsection{相关文档}


\begin{itemize}
  \item macOS 应用新手引导：\href{/start/onboarding}{新手引导}
  \item 配置参考：\href{/gateway/configuration}{Gateway 网关配置}
  \item 提供商：\href{/channels/whatsapp}{WhatsApp}、\href{/channels/telegram}{Telegram}、\href{/channels/discord}{Discord}、\href{/channels/googlechat}{Google Chat}、\href{/channels/signal}{Signal}、\href{/channels/imessage}{iMessage}
  \item Skills：\href{/tools/skills}{Skills}、\href{/tools/skills-config}{Skills 配置}
\end{itemize}



\clearpage
% === 源文件: start/onboarding-overview.md ===
\section{Onboarding Overview}


OpenClaw supports multiple onboarding paths depending on where the Gateway runs
and how you prefer to configure providers.

\subsection{Choose your onboarding path}


\begin{itemize}
  \item \textbf{CLI wizard} for macOS, Linux, and Windows (via WSL2).
  \item \textbf{macOS app} for a guided first run on Apple silicon or Intel Macs.
\end{itemize}


\subsection{CLI onboarding wizard}


Run the wizard in a terminal:

\begin{lstlisting}[language=bash]
openclaw onboard
\end{lstlisting}


Use the CLI wizard when you want full control of the Gateway, workspace,
channels, and skills. Docs:

\begin{itemize}
  \item \href{/start/wizard}{Onboarding Wizard (CLI)}
  \item \href{/cli/onboard}{\texttt{openclaw onboard} command}
\end{itemize}


\subsection{macOS app onboarding}


Use the OpenClaw app when you want a fully guided setup on macOS. Docs:

\begin{itemize}
  \item \href{/start/onboarding}{Onboarding (macOS App)}
\end{itemize}


\subsection{Custom Provider}


If you need an endpoint that is not listed, including hosted providers that
expose standard OpenAI or Anthropic APIs, choose \textbf{Custom Provider} in the
CLI wizard. You will be asked to:

\begin{itemize}
  \item Pick OpenAI-compatible, Anthropic-compatible, or \textbf{Unknown} (auto-detect).
  \item Enter a base URL and API key (if required by the provider).
  \item Provide a model ID and optional alias.
  \item Choose an Endpoint ID so multiple custom endpoints can coexist.
\end{itemize}


For detailed steps, follow the CLI onboarding docs above.


\clearpage

% === 源文件: start/wizard-cli-automation.md ===
\section{CLI Automation}


Use \texttt{--non-interactive} to automate \texttt{openclaw onboard}.

\texttt{--json} does not imply non-interactive mode. Use \texttt{--non-interactive} (and \texttt{--workspace}) for scripts.

\subsection{Baseline non-interactive example}


\begin{lstlisting}[language=bash]
openclaw onboard --non-interactive \
  --mode local \
  --auth-choice apiKey \
  --anthropic-api-key "$ANTHROPIC_API_KEY" \
  --secret-input-mode plaintext \
  --gateway-port 18789 \
  --gateway-bind loopback \
  --install-daemon \
  --daemon-runtime node \
  --skip-skills
\end{lstlisting}


Add \texttt{--json} for a machine-readable summary.

Use \texttt{--secret-input-mode ref} to store env-backed refs in auth profiles instead of plaintext values.
Interactive selection between env refs and configured provider refs (\texttt{file} or \texttt{exec}) is available in the onboarding wizard flow.

In non-interactive \texttt{ref} mode, provider env vars must be set in the process environment.
Passing inline key flags without the matching env var now fails fast.

Example:

\begin{lstlisting}[language=bash]
openclaw onboard --non-interactive \
  --mode local \
  --auth-choice openai-api-key \
  --secret-input-mode ref \
  --accept-risk
\end{lstlisting}


\subsection{Provider-specific examples}


\begin{lstlisting}[language=bash]
    openclaw onboard --non-interactive \
      --mode local \
      --auth-choice gemini-api-key \
      --gemini-api-key "$GEMINI_API_KEY" \
      --gateway-port 18789 \
      --gateway-bind loopback
\end{lstlisting}

\begin{lstlisting}[language=bash]
    openclaw onboard --non-interactive \
      --mode local \
      --auth-choice zai-api-key \
      --zai-api-key "$ZAI_API_KEY" \
      --gateway-port 18789 \
      --gateway-bind loopback
\end{lstlisting}

\begin{lstlisting}[language=bash]
    openclaw onboard --non-interactive \
      --mode local \
      --auth-choice ai-gateway-api-key \
      --ai-gateway-api-key "$AI_GATEWAY_API_KEY" \
      --gateway-port 18789 \
      --gateway-bind loopback
\end{lstlisting}

\begin{lstlisting}[language=bash]
    openclaw onboard --non-interactive \
      --mode local \
      --auth-choice cloudflare-ai-gateway-api-key \
      --cloudflare-ai-gateway-account-id "your-account-id" \
      --cloudflare-ai-gateway-gateway-id "your-gateway-id" \
      --cloudflare-ai-gateway-api-key "$CLOUDFLARE_AI_GATEWAY_API_KEY" \
      --gateway-port 18789 \
      --gateway-bind loopback
\end{lstlisting}

\begin{lstlisting}[language=bash]
    openclaw onboard --non-interactive \
      --mode local \
      --auth-choice moonshot-api-key \
      --moonshot-api-key "$MOONSHOT_API_KEY" \
      --gateway-port 18789 \
      --gateway-bind loopback
\end{lstlisting}

\begin{lstlisting}[language=bash]
    openclaw onboard --non-interactive \
      --mode local \
      --auth-choice mistral-api-key \
      --mistral-api-key "$MISTRAL_API_KEY" \
      --gateway-port 18789 \
      --gateway-bind loopback
\end{lstlisting}

\begin{lstlisting}[language=bash]
    openclaw onboard --non-interactive \
      --mode local \
      --auth-choice synthetic-api-key \
      --synthetic-api-key "$SYNTHETIC_API_KEY" \
      --gateway-port 18789 \
      --gateway-bind loopback
\end{lstlisting}

\begin{lstlisting}[language=bash]
    openclaw onboard --non-interactive \
      --mode local \
      --auth-choice opencode-zen \
      --opencode-zen-api-key "$OPENCODE_API_KEY" \
      --gateway-port 18789 \
      --gateway-bind loopback
\end{lstlisting}

\begin{lstlisting}[language=bash]
    openclaw onboard --non-interactive \
      --mode local \
      --auth-choice custom-api-key \
      --custom-base-url "https://llm.example.com/v1" \
      --custom-model-id "foo-large" \
      --custom-api-key "$CUSTOM_API_KEY" \
      --custom-provider-id "my-custom" \
      --custom-compatibility anthropic \
      --gateway-port 18789 \
      --gateway-bind loopback
\end{lstlisting}


\texttt{--custom-api-key} is optional. If omitted, onboarding checks \texttt{CUSTOM\_API\_KEY}.

Ref-mode variant:

\begin{lstlisting}[language=bash]
    export CUSTOM_API_KEY="your-key"
    openclaw onboard --non-interactive \
      --mode local \
      --auth-choice custom-api-key \
      --custom-base-url "https://llm.example.com/v1" \
      --custom-model-id "foo-large" \
      --secret-input-mode ref \
      --custom-provider-id "my-custom" \
      --custom-compatibility anthropic \
      --gateway-port 18789 \
      --gateway-bind loopback
\end{lstlisting}


In this mode, onboarding stores \texttt{apiKey} as \texttt{\{ source: "env", provider: "default", id: "CUSTOM\_API\_KEY" \}}.


\subsection{Add another agent}


Use \texttt{openclaw agents add } to create a separate agent with its own workspace,
sessions, and auth profiles. Running without \texttt{--workspace} launches the wizard.

\begin{lstlisting}[language=bash]
openclaw agents add work \
  --workspace ~/.openclaw/workspace-work \
  --model openai/gpt-5.2 \
  --bind whatsapp:biz \
  --non-interactive \
  --json
\end{lstlisting}


What it sets:

\begin{itemize}
  \item \texttt{agents.list[].name}
  \item \texttt{agents.list[].workspace}
  \item \texttt{agents.list[].agentDir}
\end{itemize}


Notes:

\begin{itemize}
  \item Default workspaces follow \texttt{\textasciitilde{}/.openclaw/workspace-}.
  \item Add \texttt{bindings} to route inbound messages (the wizard can do this).
  \item Non-interactive flags: \texttt{--model}, \texttt{--agent-dir}, \texttt{--bind}, \texttt{--non-interactive}.
\end{itemize}


\subsection{Related docs}


\begin{itemize}
  \item Onboarding hub: \href{/start/wizard}{Onboarding Wizard (CLI)}
  \item Full reference: \href{/start/wizard-cli-reference}{CLI Onboarding Reference}
  \item Command reference: \href{/cli/onboard}{\texttt{openclaw onboard}}
\end{itemize}



\clearpage

% === 源文件: start/wizard-cli-reference.md ===
\section{CLI Onboarding Reference}


This page is the full reference for \texttt{openclaw onboard}.
For the short guide, see \href{/start/wizard}{Onboarding Wizard (CLI)}.

\subsection{What the wizard does}


Local mode (default) walks you through:

\begin{itemize}
  \item Model and auth setup (OpenAI Code subscription OAuth, Anthropic API key or setup token, plus MiniMax, GLM, Moonshot, and AI Gateway options)
  \item Workspace location and bootstrap files
  \item Gateway settings (port, bind, auth, tailscale)
  \item Channels and providers (Telegram, WhatsApp, Discord, Google Chat, Mattermost plugin, Signal)
  \item Daemon install (LaunchAgent or systemd user unit)
  \item Health check
  \item Skills setup
\end{itemize}


Remote mode configures this machine to connect to a gateway elsewhere.
It does not install or modify anything on the remote host.

\subsection{Local flow details}


\begin{itemize}
  \item If \texttt{\textasciitilde{}/.openclaw/openclaw.json} exists, choose Keep, Modify, or Reset.
  \item Re-running the wizard does not wipe anything unless you explicitly choose Reset (or pass \texttt{--reset}).
  \item CLI \texttt{--reset} defaults to \texttt{config+creds+sessions}; use \texttt{--reset-scope full} to also remove workspace.
  \item If config is invalid or contains legacy keys, the wizard stops and asks you to run \texttt{openclaw doctor} before continuing.
  \item Reset uses \texttt{trash} and offers scopes:
  \item Config only
  \item Config + credentials + sessions
  \item Full reset (also removes workspace)
\end{itemize}

\begin{itemize}
  \item Full option matrix is in \href{\#auth-and-model-options}{Auth and model options}.
\end{itemize}

\begin{itemize}
  \item Default \texttt{\textasciitilde{}/.openclaw/workspace} (configurable).
  \item Seeds workspace files needed for first-run bootstrap ritual.
  \item Workspace layout: \href{/concepts/agent-workspace}{Agent workspace}.
\end{itemize}

\begin{itemize}
  \item Prompts for port, bind, auth mode, and tailscale exposure.
  \item Recommended: keep token auth enabled even for loopback so local WS clients must authenticate.
  \item Disable auth only if you fully trust every local process.
  \item Non-loopback binds still require auth.
\end{itemize}

\begin{itemize}
  \item \href{/channels/whatsapp}{WhatsApp}: optional QR login
  \item \href{/channels/telegram}{Telegram}: bot token
  \item \href{/channels/discord}{Discord}: bot token
  \item \href{/channels/googlechat}{Google Chat}: service account JSON + webhook audience
  \item \href{/channels/mattermost}{Mattermost} plugin: bot token + base URL
  \item \href{/channels/signal}{Signal}: optional \texttt{signal-cli} install + account config
  \item \href{/channels/bluebubbles}{BlueBubbles}: recommended for iMessage; server URL + password + webhook
  \item \href{/channels/imessage}{iMessage}: legacy \texttt{imsg} CLI path + DB access
  \item DM security: default is pairing. First DM sends a code; approve via
\end{itemize}

\texttt{openclaw pairing approve  } or use allowlists.
\begin{itemize}
  \item macOS: LaunchAgent
  \item Requires logged-in user session; for headless, use a custom LaunchDaemon (not shipped).
  \item Linux and Windows via WSL2: systemd user unit
  \item Wizard attempts \texttt{loginctl enable-linger } so gateway stays up after logout.
  \item May prompt for sudo (writes \texttt{/var/lib/systemd/linger}); it tries without sudo first.
  \item Runtime selection: Node (recommended; required for WhatsApp and Telegram). Bun is not recommended.
\end{itemize}

\begin{itemize}
  \item Starts gateway (if needed) and runs \texttt{openclaw health}.
  \item \texttt{openclaw status --deep} adds gateway health probes to status output.
\end{itemize}

\begin{itemize}
  \item Reads available skills and checks requirements.
  \item Lets you choose node manager: npm or pnpm (bun not recommended).
  \item Installs optional dependencies (some use Homebrew on macOS).
\end{itemize}

\begin{itemize}
  \item Summary and next steps, including iOS, Android, and macOS app options.
\end{itemize}


If no GUI is detected, the wizard prints SSH port-forward instructions for the Control UI instead of opening a browser.
If Control UI assets are missing, the wizard attempts to build them; fallback is \texttt{pnpm ui:build} (auto-installs UI deps).

\subsection{Remote mode details}


Remote mode configures this machine to connect to a gateway elsewhere.

Remote mode does not install or modify anything on the remote host.

What you set:

\begin{itemize}
  \item Remote gateway URL (\texttt{ws://...})
  \item Token if remote gateway auth is required (recommended)
\end{itemize}


\begin{itemize}
  \item If gateway is loopback-only, use SSH tunneling or a tailnet.
  \item Discovery hints:
  \item macOS: Bonjour (\texttt{dns-sd})
  \item Linux: Avahi (\texttt{avahi-browse})
\end{itemize}


\subsection{Auth and model options}


Uses \texttt{ANTHROPIC\_API\_KEY} if present or prompts for a key, then saves it for daemon use.
\begin{itemize}
  \item macOS: checks Keychain item "Claude Code-credentials"
  \item Linux and Windows: reuses \texttt{\textasciitilde{}/.claude/.credentials.json} if present
\end{itemize}


On macOS, choose "Always Allow" so launchd starts do not block.

Run \texttt{claude setup-token} on any machine, then paste the token.
You can name it; blank uses default.
If \texttt{\textasciitilde{}/.codex/auth.json} exists, the wizard can reuse it.
Browser flow; paste \texttt{code\#state}.

Sets \texttt{agents.defaults.model} to \texttt{openai-codex/gpt-5.3-codex} when model is unset or \texttt{openai/*}.

Uses \texttt{OPENAI\_API\_KEY} if present or prompts for a key, then stores the credential in auth profiles.

Sets \texttt{agents.defaults.model} to \texttt{openai/gpt-5.1-codex} when model is unset, \texttt{openai/\textit{}, or \texttt{openai-codex/}}.

Prompts for \texttt{XAI\_API\_KEY} and configures xAI as a model provider.
Prompts for \texttt{OPENCODE\_API\_KEY} (or \texttt{OPENCODE\_ZEN\_API\_KEY}).
Setup URL: \href{https://opencode.ai/auth}{opencode.ai/auth}.
Stores the key for you.
Prompts for \texttt{AI\_GATEWAY\_API\_KEY}.
More detail: \href{/providers/vercel-ai-gateway}{Vercel AI Gateway}.
Prompts for account ID, gateway ID, and \texttt{CLOUDFLARE\_AI\_GATEWAY\_API\_KEY}.
More detail: \href{/providers/cloudflare-ai-gateway}{Cloudflare AI Gateway}.
Config is auto-written.
More detail: \href{/providers/minimax}{MiniMax}.
Prompts for \texttt{SYNTHETIC\_API\_KEY}.
More detail: \href{/providers/synthetic}{Synthetic}.
Moonshot (Kimi K2) and Kimi Coding configs are auto-written.
More detail: \href{/providers/moonshot}{Moonshot AI (Kimi + Kimi Coding)}.
Works with OpenAI-compatible and Anthropic-compatible endpoints.

Interactive onboarding supports the same API key storage choices as other provider API key flows:
\begin{itemize}
  \item \textbf{Paste API key now} (plaintext)
  \item \textbf{Use secret reference} (env ref or configured provider ref, with preflight validation)
\end{itemize}


Non-interactive flags:
\begin{itemize}
  \item \texttt{--auth-choice custom-api-key}
  \item \texttt{--custom-base-url}
  \item \texttt{--custom-model-id}
  \item \texttt{--custom-api-key} (optional; falls back to \texttt{CUSTOM\_API\_KEY})
  \item \texttt{--custom-provider-id} (optional)
  \item \texttt{--custom-compatibility } (optional; default \texttt{openai})
\end{itemize}


Leaves auth unconfigured.

Model behavior:

\begin{itemize}
  \item Pick default model from detected options, or enter provider and model manually.
  \item Wizard runs a model check and warns if the configured model is unknown or missing auth.
\end{itemize}


Credential and profile paths:

\begin{itemize}
  \item OAuth credentials: \texttt{\textasciitilde{}/.openclaw/credentials/oauth.json}
  \item Auth profiles (API keys + OAuth): \texttt{\textasciitilde{}/.openclaw/agents//agent/auth-profiles.json}
\end{itemize}


API key storage mode:

\begin{itemize}
  \item Default onboarding behavior persists API keys as plaintext values in auth profiles.
  \item \texttt{--secret-input-mode ref} enables reference mode instead of plaintext key storage.
\end{itemize}

In interactive onboarding, you can choose either:
\begin{itemize}
  \item environment variable ref (for example \texttt{keyRef: \{ source: "env", provider: "default", id: "OPENAI\_API\_KEY" \}})
  \item configured provider ref (\texttt{file} or \texttt{exec}) with provider alias + id
  \item Interactive reference mode runs a fast preflight validation before saving.
  \item Env refs: validates variable name + non-empty value in the current onboarding environment.
  \item Provider refs: validates provider config and resolves the requested id.
  \item If preflight fails, onboarding shows the error and lets you retry.
  \item In non-interactive mode, \texttt{--secret-input-mode ref} is env-backed only.
  \item Set the provider env var in the onboarding process environment.
  \item Inline key flags (for example \texttt{--openai-api-key}) require that env var to be set; otherwise onboarding fails fast.
  \item For custom providers, non-interactive \texttt{ref} mode stores \texttt{models.providers..apiKey} as \texttt{\{ source: "env", provider: "default", id: "CUSTOM\_API\_KEY" \}}.
  \item In that custom-provider case, \texttt{--custom-api-key} requires \texttt{CUSTOM\_API\_KEY} to be set; otherwise onboarding fails fast.
  \item Existing plaintext setups continue to work unchanged.
\end{itemize}


Headless and server tip: complete OAuth on a machine with a browser, then copy
\texttt{\textasciitilde{}/.openclaw/credentials/oauth.json} (or \texttt{\$OPENCLAW\_STATE\_DIR/credentials/oauth.json})
to the gateway host.

\subsection{Outputs and internals}


Typical fields in \texttt{\textasciitilde{}/.openclaw/openclaw.json}:

\begin{itemize}
  \item \texttt{agents.defaults.workspace}
  \item \texttt{agents.defaults.model} / \texttt{models.providers} (if Minimax chosen)
  \item \texttt{tools.profile} (local onboarding defaults to \texttt{"messaging"} when unset; existing explicit values are preserved)
  \item \texttt{gateway.*} (mode, bind, auth, tailscale)
  \item \texttt{session.dmScope} (local onboarding defaults this to \texttt{per-channel-peer} when unset; existing explicit values are preserved)
  \item \texttt{channels.telegram.botToken}, \texttt{channels.discord.token}, \texttt{channels.signal.\textit{}, \texttt{channels.imessage.}}
  \item Channel allowlists (Slack, Discord, Matrix, Microsoft Teams) when you opt in during prompts (names resolve to IDs when possible)
  \item \texttt{skills.install.nodeManager}
  \item \texttt{wizard.lastRunAt}
  \item \texttt{wizard.lastRunVersion}
  \item \texttt{wizard.lastRunCommit}
  \item \texttt{wizard.lastRunCommand}
  \item \texttt{wizard.lastRunMode}
\end{itemize}


\texttt{openclaw agents add} writes \texttt{agents.list[]} and optional \texttt{bindings}.

WhatsApp credentials go under \texttt{\textasciitilde{}/.openclaw/credentials/whatsapp//}.
Sessions are stored under \texttt{\textasciitilde{}/.openclaw/agents//sessions/}.

Some channels are delivered as plugins. When selected during onboarding, the wizard
prompts to install the plugin (npm or local path) before channel configuration.

Gateway wizard RPC:

\begin{itemize}
  \item \texttt{wizard.start}
  \item \texttt{wizard.next}
  \item \texttt{wizard.cancel}
  \item \texttt{wizard.status}
\end{itemize}


Clients (macOS app and Control UI) can render steps without re-implementing onboarding logic.

Signal setup behavior:

\begin{itemize}
  \item Downloads the appropriate release asset
  \item Stores it under \texttt{\textasciitilde{}/.openclaw/tools/signal-cli//}
  \item Writes \texttt{channels.signal.cliPath} in config
  \item JVM builds require Java 21
  \item Native builds are used when available
  \item Windows uses WSL2 and follows Linux signal-cli flow inside WSL
\end{itemize}


\subsection{Related docs}


\begin{itemize}
  \item Onboarding hub: \href{/start/wizard}{Onboarding Wizard (CLI)}
  \item Automation and scripts: \href{/start/wizard-cli-automation}{CLI Automation}
  \item Command reference: \href{/cli/onboard}{\texttt{openclaw onboard}}
\end{itemize}



\clearpage

